% =====================================================================
%  R. Nielsen, DR. H. MARTENSEN'S DOGMATIC ELUCIDATIONS ILLUMINATED.
%  Copenhagen: C. A. Reitzel, 1850.  ENGLISH TRANSLATION, from transcription.tex (Danish) as rebuilt
%  2026-10-06 (before the first commit of the book; see its header).
%  GENERATED by ebook-method/martensens-oplysninger/tr_build.py --write from tr_template.tex and the
%  translators' parts in texts/nielsen/martensens-oplysninger/.parts/tr/ (TRANSLATOR-BRIEF.md).
%  Title: "belyste" (illuminated, examined) answers Martensen's "Oplysninger" (elucidations, literally
%  enlightenments) and the motto from Psalm 36:9; "Illuminated" keeps the play on light.
%  Conventions: TRANSLATION-PLAYBOOK.md.  Page numbers of the 1850 print in the margin.
% =====================================================================
\documentclass[12pt,a4paper,openany]{book}
\usepackage[utf8]{inputenc}
\usepackage[T1]{fontenc}
\usepackage{libertinus}
\usepackage{libertinust1math}
\usepackage{textalpha}
\usepackage[english]{babel}
\usepackage{enumitem}
\usepackage[protrusion=true,expansion=true]{microtype}
\usepackage{setspace}
\usepackage[margin=1.3in]{geometry}
\usepackage{fancyhdr}
\usepackage{hyperref}
\hypersetup{pdftitle={Dr. H. Martensen's Dogmatic Elucidations Illuminated}, pdfauthor={Rasmus Nielsen}, hidelinks}
\setlength{\headheight}{15pt}
\pagestyle{fancy}
\fancyhf{}
\fancyhead[LE,RO]{\thepage}
\fancyhead[C]{\textit{R. Nielsen: Dr. H. Martensen's Dogmatic Elucidations Illuminated (1850)}}
\renewcommand{\thefootnote}{\ifcase\value{footnote}\or *)\or **)\or ***)\else
  \arabic{footnote})\fi}
\setstretch{1.3}
\usepackage{marginnote}
\newcommand{\opage}[1]{\marginnote{\footnotesize\textit{[#1]}}}

\begin{document}

\begin{titlepage}
\begin{center}
\vspace*{3em}
{\large Dr. H. \emph{Martensen's}}\\[1em]
{\Huge Dogmatic Elucidations}\\[1.5em]
Illuminated\\[0.8em]
by\\[0.8em]
Prof. R. Nielsen.\\[2em]
\rule{8em}{0.4pt}\\[2em]
{\large Copenhagen.}\\[0.5em]
University Bookseller C. A. \emph{Reitzel.}\\[0.5em]
{\small Printed by Royal Court Printer \emph{Bianco Luno.}}\\[0.5em]
1850.\\[4em]
{\small Translated from the Danish}
\end{center}
\end{titlepage}
\thispagestyle{empty}
\vspace*{12em}
\begin{center}
``In thy light shall we see light!''\\[0.4em]
\hspace*{6em}{\small A Dogmatic Elucidation.}
\end{center}
% the motto is Psalm 36:9 (Danish: „I dit Lys skulle vi see Lys!“), given here in the Authorized Version's words
\clearpage

% --- p. 3 ---
\setcounter{footnote}{0}%
\opage{3}My ``examining review'' has thus had its effect;
Dr. Martensen, who in the Preface to his Dogmatics seemed
inclined to dispose of certain points in a rather ambiguous
and cursory manner, has thus after all at last seen the necessity
of declaring himself with somewhat greater definiteness,
and has now indeed, in his ``Dogmatic Elucidations,'' elucidated
for us a good many things which, without the ``examining review,'' would no doubt
have remained unelucidated.

That these elucidations, composed with ``so much clarity,''
must have caused a sensation in the theological reading world
and opened many eyes, I can doubt all the less,
since I too am compelled to the confession that the ``clarity''
in these dogmatic elucidations far exceeds what
I had expected in that respect. If, then, Dr. M. can for
his part be satisfied with his elucidations, then
I too, under the circumstances, am truly satisfied; indeed,
this satisfaction of mine is precisely what has brought it about that a comparatively
long time has already passed before I could again
bring myself to break off my rest and take up
the matter.

What, then, must especially be singled out as an unmistakable
% --- p. 4 ---
\setcounter{footnote}{0}%
\opage{4}merit of the ``Dogmatic Elucidations'' is the impregnable
position in which Dr. M. here places himself toward ``the critique,'' and the
evident superiority with which he asserts his dogmatic
ποῦ στῶ against his impotent ``reviewers.''

Here, naturally, there can be no thought that Dr.
M. should have been brought into any embarrassment whatever, or should have
even a single concession to make to his opponents.
On the contrary, with his dogmatic speculation everything is in the
best of order: the Dogmatics, its execution, its Introduction,
its Preface, everything here is so secure, so well grounded, so coherent,
that whoever cannot now see that the objections of ``the critique''
are as unwarranted as Dr. M.'s reviewers are
incompetent must lack all ``capacity for immediate and living
intuition of ideas.''

People have, to be sure, believed that a critique of Dr.
M.'s scientific treatment of the Christian dogmas had really
got under way; but---whether it be owing to a lack
of ``capacity for immediate and living intuition of ideas''---according
to the dogmatic elucidations this is pure imagination;
the critique has not yet begun. ``I can''---says
Dr. M. (p. 7)---``only ascribe to this whole critique a merely
provisional character, indeed must find that the critique of my
Dogmatics, insofar as such a thing could be expected at all,
cannot yet be said to have begun.''\dots{}

People have further believed that in Mr. Kierkegaard's
many writings, and especially in the pseudonym ``Johannes Climacus,''
there was a good deal that might well come into consideration
when it was a matter of assessing a dogmatician who,
after having come down considerably in speculation, has nevertheless not
been able to bring himself to give up speculation; but---
% --- p. 5 ---
\setcounter{footnote}{0}%
\opage{5}it must surely be owing to a lack of ``capacity for immediate and
living intuition of ideas''---this too is, according to the dogmatic
elucidations, a pure error. Dr. M. has (p.
13) absolutely no use for the ``doctrines'' that people have wished to communicate to him as an
extract from the pseudonym, since he
``must apply to these the old saying, that what is true in them
is not new, and what is new is not true.''

People have finally believed that the evangelical conception
of faith, as the paradox impenetrable to the natural understanding,
must, after all that recent times have
brought to light about the heterogeneity of faith and knowledge when the concept
was sharpened, nonetheless put the speculating dogmatician in some
embarrassment; but---this too must surely stem from a lack
of ``capacity for immediate and living intuition of ideas''---this
too is, according to the dogmatic elucidations, a pure
error. The knot lies in this, that one has neglected to translate
the word paradox into the mother tongue. Dr. M., who ``is
generally accustomed, when a difficult scholastic terminology
lies before him, to seek out the corresponding expression in the mother tongue,''
has then indeed, by giving several quite fitting
translations of this difficult word, discovered with his usual thoroughness
that ``the paradox,'' if one merely seeks the
corresponding expressions in the mother tongue, says absolutely nothing---which,
to be sure, can very well be reconciled with---``dogmatic
speculation.''

If the critical objections of the reviewers are really
as empty as can be seen in the dogmatic elucidations,
then neither can there be any disagreement about how the
persons who have been able to put forward such objections
are to be treated. This too Dr. M. has understood.
% --- p. 6 ---
\setcounter{footnote}{0}%
\opage{6}Instead of the usual circumlocutions by which, in literary
polemic, one otherwise seeks to expose one's opponent without
however injuring him morally, Dr. M. here avails himself, with
seriousness and authority, of ``direct communication,'' and
without beating about the bush lets the honored public know what sort of people his
opponents are.

One of these opponents, namely Mr. Stilling,
is accordingly, according to the dogmatic elucidations (cf. p.
7--8), a subject utterly miscarried in the speculative-dogmatic sense,
a person who to such a degree lacks ``positive
knowledge and an acquaintance, grounded in religious experience of life,
with the truths of revelation, that he has even
believed he could dispose, in a couple of pages, of the doctrine of the Trinity,
the doctrine of the Incarnation, the doctrine of ubiquity, the doctrine of
the Fall and original sin, of freedom and grace, besides
a whole cycle of other dogmas''; a dilettante in logic,
who is to such a degree bereft of ``capacity for immediate and living
intuition of ideas'' that Dr. M. cannot engage with him at all,
especially since he---which cannot but offend the dogmatician's
aesthetic sense---has ``unwashed hands'' and resembles
``Erasmus Montanus!'' With this ``dogmatic elucidation''
Mr. St. is thus dismissed once and for all; Dr. M. does not know
whether others will ``take closer account of the Magister's
scribbling''; His Reverence has (Oplysn. p. 8) for his
part immediately laid it aside, ``not to take it up again.''

The reviewer who, with an apparatus of proofs and reasons
such as that with which Mr. St. has equipped the above-mentioned
critique, can bring it about that the criticized author treats
him as obligingly as Dr. M. in the present case
treats Mr. Stilling, may, however, congratulate himself; for it
% --- p. 7 ---
\setcounter{footnote}{0}%
\opage{7}is surely unmistakable proof that the critique has done its work thoroughly.
Mr. St. has thus this time undeniably been luckier
than I; his critique has evidently made a more thorough
impression than mine, which is why Dr. M. also treats
him with evident partiality. Even so, I am not
envious; I gladly grant Mr. Still. all the literary
recognition he rightly deserves, and have moreover every
reason, under the circumstances, to be satisfied with the
moral testimonial that the ``Dogmatic Elucidations'' award
me; here it is, it reads thus: ``Blessed are they who
will not force their wisdom on others!'' This saying of
``the famous dogmatician, old Daub,'' Dr.
Martensen (p. 9) could not help recalling when he saw ``the
rashness'' with which Prof. N. wishes to force upon him the
wisdom which he (Prof. N.) believes he has found in Mr.
Kierkegaard's writings. Prof. N., who ``must himself admit
that he has come over to the new standpoint by a stroke of luck,''
is now moreover so vexed that Dr.
M. has not also changed standpoint, that he even wishes ``with
a vehemence bordering on fanaticism'' to give His Reverence
``an eye for the problem.'' The rash, importunate,
fanatical Prof. N. has, however, not yet (Dogm. Oplysn.
p. 11) himself become at home on his new standpoint, and
Dr. M. therefore (p. 11) cannot but find that it has
gone with him as it usually goes with beginners who adopt
a new philosophical doctrine, that they not only think with their
master's thoughts but also speak in his words, and for a time lose
the use of their own tongue. The rash, importunate,
fanatical Prof. N. is himself still (Dogm. Oplysn. p. 11)
``so entangled in the swaddling bands of the new wisdom,'' and ``
% --- p. 8 ---
\setcounter{footnote}{0}%
\opage{8}moves still so laboriously, so discipularly among the first elements of this wisdom,
that he is not even able to present to
Dr. Martensen his doctrines in his own
words.''---``That Prof. N. was not particularly strong at recognizing
specific differences and observing boundaries,'' Dr. M.
(Opl. p. 62) has long known; but now the confusion in this
rash, fanatical head goes so far that Prof. N. no longer
makes any distinction between good and evil, but without further ado
``lumps the mystery of evil together with all other mysteries, not
only supernatural but also natural.'' (Cf.
Dogm. Oplysn. p. 65); indeed---so confused is this importunate,
fanatical Prof. N., that, although evidently
seized by zeal for a one-sided ``intellectualism, criticism
and scholasticism,'' he nonetheless passes himself off
as the defender of faith (Dogm. Oplysn. p. 71), believes (Oplysn.
p. 72) that he speaks in the simplicity of faith of the incomprehensibility of divine
things, has begun to become ``absolute in faith,''
and ``has begun to speak with the wondrous parrhesia and
certainty that is astonishing to other simple believers, who
cannot so suddenly raise themselves to the highest heights of faith.''
---With all this Dr. Martensen now foresees (Oplysn.
p. 76) ``that the time will not be far off when once again
a principal change will take place in Prof. N.'s views''; for
Dr. Martensen cannot possibly suppose ``that so lively, so
vigorous, and, in spite of all its volatility and dilettantish
inconstancy, yet essentially thorough a nature can long find
itself satisfied in going the way on which it has now set out.''

See, this is, without distortion or exaggeration, the light
in which Dr. M.'s dogmatic elucidations place those
who have permitted themselves to shake the ποῦ στῶ and to doubt
% --- p. 9 ---
\setcounter{footnote}{0}%
\opage{9}the scientific character of his ``Christian Dogmatics.'' But---
is the Martensenian ποῦ στῶ now really so unshakeable,
and the security, superiority and authority with which
he appears in his dogmatic elucidations so decisive
and firmly grounded, that anyone who dares to doubt
it must without further ado put up with being regarded as a
confused, rash, importunate, fanatical head? This
question is surely hardly without all significance. We shall
now see.

Before I go further into the matter itself, however,
I must provisionally direct the reader's attention to two
points, namely: the peculiar relation in which the author
of the ``Dogmatic Elucidations'' places himself to the pseudonym,
Johannes Climacus, and the striking haste
with which he disposes of Mr. Stilling.

With great definiteness Dr. Martensen rejects any
regard to the writings published by Mr. Kierkegaard. He
thus declares (p. 12): ``As regards Johannes
Climacus and the other pseudonyms, these are in the present
context entirely irrelevant to me.'' But it
cannot be denied all the same that the apparent definiteness
with which the Right Reverend author waves the pseudonym
away (a definiteness that immediately afterwards, p. 13, is passed off
as modesty) betrays a certain inner embarrassment
which not even the dogmatic suffisance with which Dr. M.,
here as everywhere else, puts forward his assertions can conceal.
The extract from the pseudonym that makes up the first half
of my ``examining review'' undoubtedly came
somewhat inconveniently for His Rev., just because it came---from that
quarter. Dr. M. has, namely, for certain reasons most of all wished
% --- p. 10 ---
\setcounter{footnote}{0}%
\opage{10}to avoid any immediate collision with a certain
Johannes Climacus and the whole ``extensive literature''
to which this figure belongs; but when he now saw that
the basis on which my critical assessment of his dogmatic
speculation rests was an open extract from the pseudonym,
he felt, as it seems, that he could not well
engage with my ``examining review'' without saying
a couple of words about the pseudonym.

So Dr. Martensen came at last into a kind of
immediate contact with the pseudonym; but---is it
not remarkable!---His Rev. has hardly got his ``direct
communication'' begun with a couple of words before the ποῦ στῶ
trembles, and there appears an inner insecurity, an irresolution
and wobbliness that is otherwise so foreign to the Martensenian
``genius.''

The unpleasant vis à vis with the polemical pseudonym
really makes the otherwise decided dogmatician so
irresolute and vacillating that he is not even in agreement with
himself as to whether he will admit to having read the pseudonym
or not read him, understood him or not understood
him: he wants, as we shall now see, both the one and the
other, and yet nevertheless neither the one nor the other,
a state that categorically expresses vacillation.

On the one hand Dr. Martensen wants to enjoy the advantage
of his fine acquaintance with Mr. Kierkegaard's writings,
and gives it to be understood, not indistinctly, that he has a quite
different (and then naturally also a far more correct) conception
of these writings than I, who nonetheless admit to having studied
them; since he indeed (p. 12) supposes ``that these writings
move in an entirely different direction from the one taken by Prof.
% --- p. 11 ---
\setcounter{footnote}{0}%
\opage{11}Nielsen.''---On the other hand Dr. Martensen will
not lay any weight at all on his conception of the
pseudonyms, ``for''---he says (p. 13)---``my acquaintance
with this extensive literature is, as I said, only very
slight and fragmentary.''

Since Dr. M.'s acquaintance with the ``extensive literature''
of the pseudonyms is, as said, ``so slight and fragmentary,''
the extract from Johannes Climacus by which His
Rev. is so greatly embarrassed is declared with dogmatic definiteness
to be an aggregate of propositions for which no one but
I shall be responsible.

The responsibility that rests on me in this respect I
shall not disown either: I ought indeed to answer
for it that the ``propositions'' in question are the pseudonym's own,
and that they have not in anything essential been distorted by me; on the other hand
there is something else for which I dare not answer:
I dare not, namely, vouch for whether Dr. M. will admit to
having understood these ``propositions'' or not.---Since it is indeed chiefly
at the emphatic combating and thorough refutation
of these ``propositions'' that the ``Dogmatic Elucidations''
aim, one would certainly suppose that Dr. M.
would straightforwardly acknowledge having understood the ``propositions.''
But one cannot be quite sure of that after all; for on p. 17--18 His Rev. complains loudly that his
reviewer has not tried his hand at an independent exposition,
``in order''---so run his words---``to clear up these propositions, unintelligible to me
and to so many others, but\dots{}.
has fobbed us off with the pseudonym's own words, by which
we became just as wise as we were before. For if it
went thus with us with the green tree, or the pseudonym's
% --- p. 12 ---
\setcounter{footnote}{0}%
\opage{12}writings, that we could not get the right view of them, how
then shall it go with us with the dry, or with Prof.
Nielsen's excerpt?''

Thus---if I am not mistaken---Dr. M. admits that
he has not yet understood the pseudonym's ``propositions.''
Indeed he admits still more, for he admits (p. 18)
that he, together with ``so many others,'' really once
undertook to study the pseudonyms and to work his way properly
through the whole ``extensive literature,'' but has
had to give it up, whether on account of his ``individual
intellectual bent'' or some other hindrance.

Dr. M., who indeed lays so much weight precisely on
thorough research, Dr. M., who indeed precisely in his ``Dogmatic
Elucidations'' has supplied us with such valuable samples of
his learned researches in Tertullian and Anselm, the researching
Dr. Martensen surely cannot possibly complain of not
being able to understand an original author when he will not take
the trouble to research,---``\textit{negligentiæ mihi videtur};
yes, \textit{negligentiæ mihi videtur}!'' Dr. M. has thus in earnest
researched into the pseudonyms, researched after meaning and ``coherent
cognition,'' but---his earnest research
has been in vain.

Yet, what am I saying, Dr. M. has of course not researched
into the pseudonyms at all, not even taken the trouble to
read them properly through; he confesses it himself: ``my
acquaintance with this extensive literature is, as I said, only
very slight and fragmentary.''---\textit{negligentiæ mihi videtur},
yes, \textit{negligentiæ mihi videtur}!---but he confesses it
himself, and one must take so earnest a man at his word.

Dr. M. can thus declare the pseudonyms (``the green
% --- p. 13 ---
\setcounter{footnote}{0}%
\opage{13}tree'') unintelligible, although he has not really
read them at all; Dr. M. can thus very well refute the pseudonym's
``propositions,'' but he cannot understand them: see, these are
points for which I will not be responsible\footnote[1]{Dr. Martensen complains that he cannot understand the pseudonyms, but can find in these writings only what in the Preface to his Dogmatics he calls ``stray thoughts, aphorisms, fancies and flashes.'' This complaint is rejected as unwarranted. When one reads an author only cursorily and ``fragmentarily,'' how should one then be able to find more than stray thoughts and aphorisms? Dr. M.'s ``direct communication'' that he has not been able to understand the pseudonyms is thus here to be taken as ``an indirect communication,'' namely that His Rev. has neglected to read.}.

Yet about the pseudonym's ``indirect communication,'' as well
as about Dr. Martensen's and my relation to it, I shall
express myself somewhat more fully later on; I now pass to
the second point: the treatment of Mr. Stilling.

Among the grievances that Dr. M. advances against my
``examining review,'' that one is especially emphasized according to which I am
supposed not to have found myself able to render the pseudonyms'
problem in my own words, or to give a proper account
of how the alleged discord between faith and knowledge
actually arises; but instead to have, in a quite discipular
manner, confounded indirect and direct communication
with one another.

If Dr. M. now really in earnest lays any
weight on this complaint brought against me, if it is now
really his sincere wish to have the genesis of the problem
described, paraphrased and made clear in a direct
and didactic form: must one not then wonder greatly
that His Rev. has been able, with so extreme, or rather, with so inward
an indignation, to reject Mr. Stilling's
% --- p. 14 ---
\setcounter{footnote}{0}%
\opage{14}contribution? Everything that the Right Reverend dogmatician so greatly
misses in me he will surely be able to find precisely in Mr.
Stilling.

When Dr. M., then, dwells on my great incompetence,
namely that I am ``so entangled in the swaddling bands of the
new wisdom, and still move so laboriously, so
discipularly among the first elements of this wisdom, that I am not
even able to present my doctrines to him in
my own words''; when Dr. M., I say, from the exalted
ποῦ στῶ of his old wisdom looks down on my lowliness and
surveys my incapacity: then His Rev. ought at least
to grant Mr. Stilling a milder, more inviting
glance, and rejoice that the Magister is in that respect
far more forward than I.

For Mr. Stilling is not, like me, so entangled
in the swaddling bands of the new wisdom that he should have
lost the use of his own tongue, for he really is able
to present to Dr. Martensen ``his doctrines in his own % print: Læresætuinger (misprint)
words''; Mr. Stilling likewise does not have, like me, the
misfortune of confounding the indirect communication with the direct,
for his communication is as direct and didactic
as possible. Since, however, I see to my great astonishment
how Dr. M. has laid Mr. St.'s work aside,
not to take it up again, I will, to prove what I
say, on this occasion take up this work again, and in a few
main features remind the reader of its contents.

In his treatise, worked out with logical acuity, Mr. St. treats
precisely the two ``doctrines'' on which everything
in this investigation especially depends, namely I) \textit{Quæro intelligere},
\textit{ut credam}, and II) \textit{Credo, ut intelligam}. Thus:

% --- p. 15 ---
\setcounter{footnote}{0}%
\opage{15}``I) \textit{Quæro intelligere, ut credam} (I seek, through
appropriation in thought, through scientific cognition, scientific
demonstration, to revive my faith, somewhat enfeebled by reflection).''

Under this heading Mr. Stilling (``On
the Imagined Reconciliation of Faith and Knowledge,'' p. 3--44) treats the
many-sided relation in which human thinking can
place itself to its object, and then proves from this that the religious
object, the object of faith and hope, differs
from every other object of thought in this, that when the tension
reaches its highest, it even repels and frightens back the subjectivity
striving after knowledge. The deeper nature of this repelling
relation is, however, so far from becoming clear to the reflecting
person at first glance that, on the contrary, the human spirit working
in theories has had to undergo
many and painful experiences before it learned to see
what this means. For thinking fertilized by faith actually begins by regarding the Christian object
like any other object, which must nonetheless gradually
allow itself more and more to be worked up and appropriated in scientific,
objective knowledge.

``Let the Christian even''---thus the author,
p. 7--8, has the spirit speak in monologue---``call itself `a
hard food'; my power of thought, which is furnished with such strong
digestive organs, will surely digest it; indeed the task of
thinking it shall and must succeed; for only thus is it
really mine, only thus is it really made inward,
only thus am I, over against it, a free infinite subjectivity,
freed from the intolerable pressure of all external authority;
only thus is my relation to the Christian object a
% --- p. 16 ---
\setcounter{footnote}{0}%
\opage{16}relation that is worthy of the human spirit, for---`der
Mensch, da er Geist ist, darf und soll sich selbst des Höchsten
würdig achten', von der Größe und Macht seines Geistes
kann er nicht groß genug denken; und mit diesem Glauben
wird nichs so spröde und hart seyn, das sich ihm nicht eröffnete.
Das zuerst verborgene und verschlossene Wesen des
Universums hat keine Kraft, die dem Muthe des Erkennens
Widerstand leisten könnte; es muß sich vor ihm aufthun,
und seinen Reichthum und seine Tiefen ihm vor Augen legen
und zum Genusse geben.''

It is Hegel's spirit that speaks; it is the distinguishing mark of theoretical
knowledge that it always presupposes a direct
relation between the cognizing subject and the object of
cognition, a relation in which subjectivity must in the end become
the overreaching. On this point there is thus only the
difference between the Hegelian and the theological speculation
that Hegelian thinking does not stop halfway
but carries the consequence through. Only through this bold attempt
did it then become fully manifest that the Christian object,
the object of faith and hope, is heterogeneous with that of science,
and that there is an infinite difference between
resting in objective certainty and walking in faith and in hope.
The attempt proved to be (work cited, p. 12) ``an overreach, against
which the Godhead, as it were, protests with the words: `no!
here it is not you, man, but I, the Lord, who am ``the overreaching,''
who goes beyond your concept, a cross for your
thought, a cross that can be borne by your head only by
taking every thought captive under the obedience of faith.'\,''

This ``obedience of faith'' is, however, not to be confused % print: imidleitid (misprint)
with the insinuating submissiveness with which a mystically
% --- p. 17 ---
\setcounter{footnote}{0}%
\opage{17}unclear thinker may take it into his head to let (p. 14) ``his power of thought
be engaged as the yes-man of the dogmas; an engagement
that is no doubt sometimes offered it by dogmatizing thinkers.
For just as a yes-man is always engaged according to the method
that he apparently, and for the sake of appearances, retains
a free, independent opinion, while in reality he is
only to say what his master wants to hear: so one may
also sometimes hit upon wishing to engage the power of thought
on the same degrading conditions, and thereby obtain
its servile submissiveness, a submissiveness and a relation
which is indeed entirely different from the relation
in which thought, after a knightly struggle, at last surrenders itself
captive under the obedience of faith.''

Under the continued effort to appropriate the thoughts of
revelation, ``manifold disappointments
and nuances in the reflecting subject's relation
to the object'' may now no doubt show themselves; ``but''---it says
p. 15---``essentially only three forms are possible in this relation,
according as either I. predominantly the independence of the object (the so-called thoughts of revelation)
is asserted; or II. predominantly
the independence of the thinking subject is asserted; or III. both
the independence of the object and that of the subject are asserted at the same time.''
Under I. Mr. St. classes the sort of thinkers who undertake
``to think and to present as thinkable the whole
so-called old orthodox Christianity'' (Göschel).

Standpoint II. is represented by Hegel and Marheineke.
On this standpoint ``the tension between object
and subject has really fallen away.'' The Christian object
was here really reduced to an unstamped, unresisting material,
which had to put up with being corrected according to the demand of thought.
% --- p. 18 ---
\setcounter{footnote}{0}%
\opage{18}``What (p. 20) was needed above all was a
critique that put an end to the sham union (which had only been
achieved by the subject's overreach over the object), and that
presented the Christian in its sharp contours, so that
it must become truly conspicuous to the subject that the object
(the dogma) would be a hard food for its thought.''

Only at standpoint III. are the sides even. ``The object
is in this form of the relation more than material, the subject
more than yes-man. And, since the independence of both terms
is thus asserted at the same time, it must be possible for it to come to
a decision whether the independent subject is able, through a process of cognition,
to appropriate the undistorted object or
not.''---How ``the tension'' arises here, and what
it essentially signifies, are now the cardinal points that Mr.
St. treats with as much warmth as penetrating
acuity.

Mr. St. conceives ``the tension'' from two sides: the
theoretical and the practical. ``In the theoretically tense relation
(p. 23) the intensity and extent of the tension depend
on the intensity and extent of the subject's power of reflection.
The more intense the power of reflection, the stronger the theoretical tension;
and it is therefore by no means every subject's affair
to come into this tension. In the personally tense relation,
on the other hand, essentially every person, even the simplest, stands,
as soon as he is so chosen by existence for
tribulation and pain that what---seen and conceived in a
human way---must be felt by him, the person concerned,
and conceived as the deepest misfortune and adversity, which crushes
and hampers his life, is nonetheless, seen and conceived in a Christian
way, supposed to be the best, God's governance, good fortune. But
% --- p. 19 ---
\setcounter{footnote}{0}%
\opage{19}only the union of the theoretically tense and the practically
(personally) tense relation in one subject forms the highest point of the tension
and the life-situation in which the reflecting
subject must be able to learn a little about what it
is that is really demanded of him when he is bidden to believe.''---
The reader who, with sound earnestness and quiet reflection, will
test what the treatise (p. 24--44) communicates about the
Christian faith in providence and the Christian relation of prayer shall in
truth ``get an eye'' for the genesis of the problem and discover
a ``paradox'' that by no means lets itself be adapted to dogmatic
speculation by one's seeking ``the corresponding expression
in the mother tongue.''---``II) \textit{Credo, ut intelligam} (I, a reflecting person, have
fortunately escaped from the paths of reflection back to a firm,
unshaken biblical-ecclesiastical \textit{Credo}, and now seek, starting
from it, once again to raise more and more my certainty of faith to the
certainty of intelligence).''

Under this heading, then, Mr. Stilling treats
(p. 45--92) the proposition which may well rightly be called
Dr. Martensen's cardinal proposition, the proposition which really
designates our speculative dogmatician's ποῦ στῶ; from which
it can provisionally be seen that Dr. M. does his reviewers
a great injustice when he accuses them of not wanting
to acquaint themselves at all with his ``fundamental view,'' and begins his
``Dogmatic Elucidations'' by hurling Luther's words:
``Ihr habt einen andern Geist als wir!''---against them.

Mr. St. has already introduced the detailed discussion of ``\textit{Credo}, ut \textit{intelligam}''
in his preceding section; for
he points out (p. 37) that it goes, namely, no
``better with reflection \textit{post fidem} (i.e. after I have
% --- p. 20 ---
\setcounter{footnote}{0}%
\opage{20}taken the side of faith) than \textit{ante fidem} (i.e. before I took the side of faith,
but was still stuck right in the middle of the oppositions of reflection).
The opposition of reflection (unbelief) was overcome at the first
step of faith, not by reflection, but \emph{by} and
in believing life, which dictated obedience.'' The same
will and must now happen at each of the following steps.
Therefore also (p. 38) ``certain dogmatizing thinkers'
distinction between a speculative proof \textit{ante fidem}
(which is supposed to be impossible) and a speculative proof \textit{post fidem}
(which is supposed to be possible) is a completely false and untenable
distinction.''---Faith is, in other words, quality, and
its personal conviction qualitatively different from reflection,
which always has not only a testimony \textit{pro}, but
also a testimony \textit{contra}.

If now Dr. Martensen's Dogmatics were really soundly
built and scientifically grounded on the proposition:
\textit{Credo, ut intelligam}, then, according to Mr. Stilling's
contention, it would have to give, both in its general design and in the treatment
of the individual dogmas, a clear and complete
testimony to how faith (credo) and reflection
(\textit{ut intelligam}) actually relate to each other. The dogmatician
ought thus to show in deed that he knows both
the quality of faith and the \textit{genus cognoscendi} that faith
requires, and also knows the quality of objective knowledge (intelligence)
and the result to which a free and many-sided
reflection must necessarily lead. But precisely in this respect
Mr. Stilling thinks that Dr. Martensen's ``coherent cognition''
does not quite hang together.

Mr. Stilling objects that one cannot see
what it is that Dr. Martensen, on his wanderings of reflection, from
% --- p. 21 ---
\setcounter{footnote}{0}%
\opage{21}faith to intelligence and from intelligence to faith has actually
experienced.

``Did he experience''---asks the reviewer (p. 50)---``that
he essentially and principally could not think the Christian?''

This question is answered in the negative. ``True''---it says
p. 53---``Prof. M. says (Dogm. p. 108) that his
attempted knowledge, his speculative cognition, is not entirely
adequate, that in the highest speculative knowledge there is a moment
of non-knowledge. But in the first place this expression is
only an enfeebled and highly incorrect expression for the cross of thought,
for the offense, the paradox and the testimony of reflection \textit{contra};
and secondly, if the expression `inadequate knowledge' on closer
explanation comes to mean nothing other than a
knowledge which, qua knowledge, really does not deserve the name of knowledge
at all: then the whole result of the wandering
on `the way of human cognition,' the whole result of
the attempted knowledge, really becomes an attempted non-knowledge that worships God's
hidden wisdom. But wholly and fully and unreservedly
to admit a non-knowledge, or better (for this is the
most correct formula), that by reflecting to the bottom one reflects
oneself into the deepest and most terrible oppositions of reflection
of \textit{pro} and \textit{contra}, and stumbles on the cross of thought: that % print: Tankeforset (misprint)
is what our dogmatician balks and shrinks from.''

``Did he then experience''---the reviewer asks further (p.
62)---``that he essentially and principally could think
the Christian?''

This question too must be answered in the negative.
``It will in this test''---says the reviewer (p. 66--67)---``become quite evident that our dogmatician's
% --- p. 22 ---
\setcounter{footnote}{0}%
\opage{22}`\textit{Credo}' is in practice a scholastically inhibiting \textit{Credo},
which places the freedom of reflection under censorship and permits
it no utterances other than such as are pleasing to the dogmatician's
ear; a most peculiar way of proceeding
in order to question and sound out the opinion of the authority of thought
concerning the dogma.''

To show more closely, now, how Dr. Martensen
has made use of his credo, \textit{ut intelligam}, Mr.
Stilling singles out various dogmas, and each time comes
to the result that the dogmatician, who ``will not remain standing
at \textit{Ita} but advance to \textit{Quare},'' does not on any point
keep what he promises, but with each dogma pretends that
he is thinking, and yet does not think.

This is seen, first, in the treatment of the dogma of the Trinity
(p. 68--81). Already at the outset Dr.
M. commits the inconsistency of wanting to prove the threefoldness in God's
ontological self-revelation without, however, wanting to prove God's
existence. ``To prove''---it says, Dogmat. p. 89---``the
existence of the revealed God cannot be the task.''
To which Mr. St. remarks: ``If it can be the task of
the dogmatician to prove that the personality of the revealed God
must be thought as a threefold personality, why then can
it not be a task to prove the existence of the revealed God,
or to prove his personality? For to prove his
existence is precisely to prove his personality.''---But neither
does the test which Dr. M. here gives of his rendering comprehensible
of God's Trinity contain the least contribution toward
carrying this task further along the way of speculation. Dr. Martensen
does indeed say (Dogm. p. 128): ``To think the essential Trinity
ontologically is to think the necessary fundamental form
% --- p. 23 ---
\setcounter{footnote}{0}%
\opage{23}of God's personal life, to think the moments in God's
essence without which personality and self-consciousness are unthinkable.''
But Mr. Stilling's critique makes it only all too evident
that the Right Reverend dogmatician has not at all
thought anything of what really should be thought here. The well-known
``moments in God's essence,'' which were to supply the foundation of the concept of the Trinity,
are borrowed from thinkers whose
names Dr. M. has not even mentioned. This concept of the Trinity,
borrowed from speculation, shows itself here still just as
heterogeneous with the Christian doctrine of the Trinity as in earlier
speculative dogmatics; which is why the old
problem, whether there must be thought to be one or three
I's in the Godhead, is not solved at all in the Martensenian
``speculation,'' which merely helps itself by hushing up reflection
and passing over the difficulties in silence.

In similar fashion Mr. Stilling continues his critique
of the speculative thoroughness with which Dr. Martensen
treats the remaining dogmas, namely the dogma of the Incarnation
(work cited, p. 81--85), the dogma of original sin and personal
accountability (p. 81--88), the dogma of immortality
(p. 88--92). That the critique of these individual dogmas is
somewhat brief must no doubt be explained by the fact that the honored
reviewer has not had ``immediate and living intuition of ideas''
enough to dwell properly on the positive apparatus of
scriptural passages or on the ingenious author's witty paraphrases
of what has been handed down, but has instead, without ceremony,
laid bare the real nodal points, in order thereby
to establish that Dr. M. has at bottom thought nothing at all
of what nonetheless essentially ought to be thought here.

Whether, for that matter, there is anything much to
% --- p. 24 ---
\setcounter{footnote}{0}%
\opage{24}the critique sketched here I will, however, rather leave to
the reader himself to consider; but I cannot, on the other hand, refrain
from expressing my astonishment at Dr. Martensen's striking
aversion to the Stillingian contribution.

Dr. Martensen will now, as said, know nothing of Mr.
Stilling's critique; but, in my humble opinion, His
Rev. ought from his point of view precisely to have laid considerable weight
on the Stillingian critique, and, when he found that he
ought in any case to engage with only one, then rather
to have chosen Mr. Stilling than me; His Rev. ought,
if I am not mistaken, evidently to have done so out of two considerations:
partly for his own sake, partly for people's sake.

As regards the first consideration, Dr. M. indeed complains
that he cannot properly understand ``the indirect communication''; all the
more, then, ought he to have appreciated that Mr. St.'s
work, by its ``direct communication,'' would be precisely
a good guide for him to get properly into the matter, and
see what it was that was here being passed off as ``the problem.'' For
the problem lies, as Mr. St. shows, quite rightly in this,
that the object of faith stands in an entirely different relation to
human cognition than any whatever of the objects
that otherwise become the object of objective knowledge.

Had Dr. M. now fixed his attention on this
\emph{inverse} relation of the object of faith to the knowing subject,
had Dr. M. engaged with this difficulty instead of
helping himself with pure assurances and ``further'' continuing
(Dogm. Oplysn. p. 55) ``to regard the dogmatician's
relation to revelation in a manner corresponding to that of the natural scientist
to nature''; had Dr. Martensen shown point for
point the emptiness of Mr. Stilling's arguments
% --- p. 25 ---
\setcounter{footnote}{0}%
\opage{25}(especially the one with the ``yes-man'' and ``the series \textit{contra}''); indeed, had
Dr. M. really felt disposed to parry the Stillingian
attack with the security, calm and superiority which
the consciousness of a just cause gives: much would have been
gained here. Not only would Dr. M., as said, by this
procedure have got (which he now unfortunately has not)
properly into the matter; but he might then perhaps also have
succeeded even better in ``orienting'' his readers ``about
the real state of the matter.''

Indeed, had Dr. M. merely been able to condescend to refute
Mr. Stilling's critical objections: much would certainly
have been gained here.

All well-disposed readers (especially all those readers who are so
well-disposed as to content themselves with Dr. M.'s ``Dogmatic Elucidations''
without troubling themselves much about the opponents' critical ``scribblings'')
would then no doubt have reasoned thus: ``See,
here one sees what the new paradox-mongering amounts to;
the pseudonym is indeed only a humorist, a `gymnast.' With
the humorous gymnast so earnest a dogmatician
as Dr. M. cannot engage, since he naturally cannot
know whether the humorist's `indirect communication' is to be
jest or earnest. Prof. N. is admittedly no humorist,
but still an imitator as frivolous as he is unfortunate, who
has become infatuated with the humorist and has turned rash.
With the `rash,' `importunate' and `frivolous'
Prof. N. so judicious and thorough a researcher as Dr. M.
cannot really engage either; all the less since Prof.
N. only gives us an `excerpt' from the pseudonym, and is not able
to develop the problem in `direct communication.' But see,
hardly does an opponent appear on whose didactic
% --- p. 26 ---
\setcounter{footnote}{0}%
\opage{26}attempt one can after all to some extent reflect, hardly does a critic present himself
who is after all to some degree able to use his
own tongue, hardly does Mr. Stilling appear with his direct
communication and his logical explanatory grounds, before it
also turns out that in Dr. Martensen we have the man who
likewise, in direct communication and with still weightier explanatory grounds,
is able to `sweep before the door of dogmatics'
and `orient us readers about the real state of the matter.'\,''

When Dr. M., notwithstanding this, lets this evident
advantage slip through his hands and, although he so strongly
calls for ``direct communication,'' nevertheless declines any
negotiation with the opponent who offers him precisely
what he wishes: then one is not, after all, so entirely
``oriented about the real state of the matter.''

The self-thinking reader (the reader who does not dare
without further ado to regard Dr. M.'s Elucidations as articles of faith)
must then involuntarily ask himself: ``What can it be
that has moved the Right Reverend dogmatician to
treat Mr. St. in so scornful a manner?''; for, truth
to tell, here Dr. M. speaks by no means as the man of God
who thinks only for ``the holy catholic Church,'' but
rather as the deeply offended author who is led by an all too
natural self-regard into an imprudence: \textit{tantæne}
\textit{animis coelestibus iræ}! What sin has Mr. St. then committed, that he
was so abruptly cast off by a friend of ``the congregation,'' by the Right Reverend servant of the Word? The transgression can surely not consist in personal rudeness and moral invective;
for the critique in question contains no utterance that anyone may not
honorably permit himself in such a matter, provided only that his arguments otherwise
hold. But in what, then, does Mr. St.'s
% --- p. 27 ---
\setcounter{footnote}{0}%
\opage{27}misdeed consist? It surely could not consist in his
having spoken so freely about reflection as ``yes-man'' and ``the series
\textit{contra}'' that Dr. M., who has no need at all of ``the series \textit{contra},''
since he speculates solely in regeneration, has
concluded from it that ``this unworthy reviewer was not yet regenerate''?
The misdeed surely could not consist in Mr.
St.\footnote[1]{A Couple of Questions to Prof. Scharling.} having been so immodest as publicly to mention a fact
which is of so delicate a nature that Dr. M.'s ``modesty''\footnote[2]{Dogm. Oplysning p. 69.}
only permits him to regard it as ``less essential,''
the fact, namely, that the Right Reverend dogmatician
had, according to his own express, printed and published
testimony, an entirely different concept of science
in the year 1839 than the one that now, from the year 1849, lies
at the basis of his Dogmatics. Should this be the case,
then His Rev. has surely again, against his will, come
to give us---``an indirect communication.''

In order at once to ``orient'' us about ``the real state
of the matter,'' Dr. M. declares, as we have seen, that the critique
of his Dogmatics cannot really be said to have
begun yet. Now I can very well understand what sort
of lighting effect the Right Reverend author intends by this
``dogmatic elucidation,'' but the significance of the grounds
on which he supports his assertion I cannot, on the other hand, properly
understand. He complains that the critique has kept chiefly
to the Introduction. But even if this were the case,
this circumstance would still here ground no
complaint. Indeed, if the critique aimed at showing that the work's form
% --- p. 28 ---
\setcounter{footnote}{0}%
\opage{28}and execution did not correspond to the principles set up in the Introduction,
that would be another matter. But here the matter
stands the other way round. The critique concedes, on the contrary,
that the execution is as good as one can in fairness demand it to be according to the stated
``principles''\footnote[1]{``There is, to be sure, in this Christian Dogmatics, seen from the side of form, also something pleasant, something appealing, something that in certain moods acts, if not exactly rigorously convincing, then still gently reassuring. There is in this Christian Dogmatics something many-sided, something reconciling, something peaceable, something pliant. There is in this Christian Dogmatics, when it is regarded as a systematic whole, also a certain something artistic, something architectonic, something that at one and the same time as it were weds the University building with the Palace Chapel, science with religion: sharp understanding and naive simplicity, modern reflection and mystical immediacy, the clarity of the concept and the obscurity of mystery, all melts together in the unity of style, in festive calm, solemn twilight, like the gleam of day with the glow of an altar candle.'' Undersøg. Anm. p. 43.}; only
the scientific tenability of the principles is what the critique
contradicts. Since Dr. M., as he himself also (p. 5.) confesses in the ``Dogmatic Elucidations,''
has set forth his principles in the Introduction
``as sharply and definitely as
possible,'' I really do not see what can be objected
to a critique of the ``principles'' keeping chiefly
to the ``Introduction,'' and only partly fetching
examples from the execution.

Further it is said (p. 6) that the critique is now again supposed
to have come onto ``the old, broad road which in a bad
sense may be designated as the German one, where one namely
tumbles about in nothing but propaedeutic investigations, raises
clouds of dust, but does not get to the matter at all.''---But
should it now really be entirely superfluous to go
into ``propaedeutic investigations,'' and to return to
% --- p. 29 ---
\setcounter{footnote}{0}%
\opage{29}the principles, when one wishes to find out whether a dogmatic
work has any principle or is---unprincipled?

At the time when Dr. M. wrote his little work ``On the Autonomy
of Self-Consciousness,'' and ``in a few pages'' criticized
Kant and Schleiermacher and Hegel alike, he was no
doubt himself of the opinion that one might well test
a dogmatic system without exactly losing oneself in the detail of
the dogmatic material; at that time he did not yet regard it
as superfluous to dwell on the introductory questions
and the ``principles''; but now, when his own dogmatic
work is completed, now it is different; now nothing at all
is to come of discussing the Introduction to the Dogmatics,
now one at once becomes an ``Erasmus Montanus'' when one
has the ``importunity'' to want to discuss the dogmatic
principles. A significant change!\footnote[1]{Dr. Martensen, who made his fortune at the University precisely by lecturing on the philosophy of standpoints in a brief, popular-aesthetic survey; Dr. M., who indeed rose precisely through his ``introductions on faith and knowledge'' and his ``preliminary arrangements for the construction of a system''; Dr. M., who at the beginning announced himself as a man of the idea who would touch only the nerve of systems: this same Dr. M. has now (Dogm. Oplys. p. 5) all at once become so sated and weary of all ``theological and philosophical introductory questions'' that he has even begun to transform himself into a man of material, who talks of learned researches in Tertullian. If only it should not befall Dr. M. that over his deep research in Tertullian he forgets the Acts of the Apostles; for to confuse Festus and Agrippa (cf. Dogm. Oplys. p. 30 with Acts ch. 26, 24) is---\textit{NB.} for a learned researcher---a little unfortunate.} % print: ubeldigt (misprint)

To want, as Dr. M. (Dogmatisk. Oplysn. p. 7.)
expresses it, ``to dispose of the doctrine of the Trinity, the doctrine of the Incarnation,
the doctrine of the Fall and original sin, of freedom
and grace, besides a whole cycle of other dogmas in
% --- p. 30 ---
\setcounter{footnote}{0}%
\opage{30}a couple of pages,'' is no doubt unreasonable when a
speculative-dogmatic exposition of each individual dogma by
itself is demanded. But this much Dr. M. will surely still, when he
properly bethinks himself, readily concede to us, that there is a great
difference between supplying a positive elaboration of a
dogma and subjecting such an elaboration to a
scientific critique. The former will as a rule require many
pages; the latter can, if only the critic hits the right
point, very well be done in a couple of pages.'' In this
respect too it would have been desirable if Dr. M. had not merely
leafed through Mr. St.'s treatise and counted its pages, but
had also read and thought through what stood on these pages,
before he so ``immediately laid the Magister's scribbling aside, not
to take it up again.''

As proof that there is something in what I say,
I shall merely permit myself to cite one single passage from
the Stillingian book. With regard to the dogma of ubiquity,
which is treated on p. 83, it says in a note on p. 73:
``That Prof. Martensen in reality is not able even in the remotest
degree to think the omnipresence of a punctual personality,
he also betrays quite naively in the doctrine of ubiquity,
where the talk is of the ascended Son of God, who
has his seat at the right hand of the Father, which is everywhere.
`If we would'---thus says our dogmatician, p. 387---`hold
fast and carry through the notion of an unlimited
presence of Christ in heaven and on earth, then we could
not avoid volatilizing Christ's individuality' (indeed
it is straightforwardly admitted that an unlimited presence of Christ
leads to a pantheistic conception); `for even a glorified
individuality, even a spiritual body, cannot be
% --- p. 31 ---
\setcounter{footnote}{0}%
\opage{31}thought without limitation.'\,'' Here the reviewer now asks: ``Can
it then in the least degree be better \emph{thought} that Christ
is personally present at each Lord's Supper, when in one town,
or in 10 to 100 different towns, the Supper is held simultaneously
at 10 to 100 different places in each town?''---

This little note, which does not take up one page, let
alone ``a couple of pages,'' is of such a nature that it not
only illuminates the Martensenian treatment of the dogma of ubiquity,
but also the principled relation in which Dr. M. as a
speculative thinker has placed himself to the ecclesiastical dogmas
as a whole.

A scientific review of our speculative-Christian
Dogmatics can, on the basis of this note, be briefly drawn up
thus: ``An original speculative dogmatician has
arisen among us. That this dogmatician is both original
and speculative is seen from his treatment of all dogmas
without exception; but that he is really as original
as speculative, and as speculative as original, appears
in particular from his treatment of the doctrine of ubiquity.
This dogmatic problem neither the older nor the
younger theologians have hitherto been able to solve in any really satisfactory
way. Only in the present work has
`the dogmatic genius' at last succeeded in solving the problem and
sublating `the one-sidedness' by an entirely new discovery. The discovery
is this: in order to be present in the Supper,
the glorified body of Christ does not yet really need to
be \emph{omnipresent}; it need only be
\emph{multipresent} at once. Now since such
a multipresence must really be \emph{thought} as a
merely relative and limited omnipresence, there is surely
% --- p. 32 ---
\setcounter{footnote}{0}%
\opage{32}nothing in the way of such a limited omnipresence
being very well reconcilable with the necessary limitation of the glorified individuality.
See, that is the solution of the problem,
a solution in which there is indeed no paradox!---How
we others, who have not been able to make the discovery, are meanwhile
to go about getting the newly discovered
`multipresence' really thought, our dogmatician does
not say: he says only that it can be thought and ought to be thought and
must be thought; but more he does not say, and this is precisely the
original. How `the genius himself' goes about it when
he is to analyze logically the concept of `relative omnipresence,'
he does not tell us either; he merely says credo, ut
\textit{intelligam}; but as soon as he merely says: credo, \textit{ut intelligam},
it is also as if it had really been thought, and
this is precisely the speculative.''

Should it still be necessary, in order to characterize the actual
standpoint, to go into ``the
whole fullness of dogmas,'' or can we perhaps let it
rest with this meager review, and say: sat sapienti?

``If the dogmatic genius''---says Dr. M.
(Dogm. Opl. p. 54)---``was productive then, it may
well be productive still,'' and in this I must grant him he is
right. But it does not now depend solely on
``a dogmatic genius'' being ``productive''; it surely also depends
essentially on what it is that ``the dogmatic genius''
produces.

Dogmatic productivity can surely not, at the present time,
consist chiefly in discovering new articles of faith
or bringing forth new dogmas. I for my part
cannot find at all that Dr. M.'s ``dogmatic genius''
% --- p. 33 ---
\setcounter{footnote}{0}%
\opage{33}has produced a single new dogma (unless it should
be the dogma of the ``multipresence'' of the glorified Christ).
The aggregate of articles of faith that he treats
in his Dogmatics is precisely the same cycle of
articles, scriptural passages and questions that have long been worked over
in the most diverse ways in the Lutheran Church.
The appropriation of this material requires neither any excessive
store of ``positive knowledge'' nor any peculiar
``acquaintance, grounded in religious experience of life, with the truths of
revelation''; one finds the whole apparatus in almost
every dogmatic compendium, indeed in some handbooks (e.g.
in Bretschneider's Dogmatics and Hase's Hutterus redivivus)
even far more fully than in Dr. M.'s book. The
Martensenian productivity has, then, in this respect no learned
significance, as though his Dogmatics were a new and distinctive
result of research in the history of dogma. The significance
of Dr. M.'s ``productivity'' must thus here be sought in
something quite different, namely in the scientific competence
with which he has known throughout how to clarify the relation
between faith and cognition.

In this respect, then, Dr. M. demands that, in order
properly to assess his ``dogmatic productivity,'' one shall (Dogm.
Oplysn. p. 6) ``let the judgment of the theory of cognition go
hand in hand with an examination of the development of the dogmas itself,
since these two mutually illuminate each other, and
cannot rightly be understood except in combination.''---I shall
now, supporting myself on Mr. Stilling's note, at once
have the honor of complying with this demand.

``The doctrine of Christ's seat at the right hand of the Father''---it says,
Dogm. §~177---``is developed in the doctrine of Christ
% --- p. 34 ---
\setcounter{footnote}{0}%
\opage{34}as the one who, by an all-pervading presence, fills
all in all (τοῦ τὰ πάντα ἐν πᾶσι πληρουμένου).''---Here,
then, is a single, definite, positive dogma. To this dogma
cognition can now place itself in a twofold relation, a scientific
and a religious, an objective and a subjective relation.
In both these relations one can begin with the proposition
\textit{Credo, ut intelligam}; but, as we shall at once see, the
proposition takes on an entirely different meaning in religious cognition
than in scientific.

In the objectively scientific sense \textit{Credo} aims at
a historical acceptance of a given fact: \textit{Credo}, I believe,
i.e. I accept as decidedly certain that the dogma in its definite
form expresses the true and right meaning of Christianity and of the Church.
\textit{Credo}, I accept the dogma as it is, \textit{ut intelligam},
i.e. in order that I may thereby come to cognition.
But to which cognition? To a scientific cognition,
i.e. to cognition of whether or not I can, by the way of science,
come to see the objective truth of the dogma.
To decide this, I must at once turn
my attention to two main points, since I must namely
investigate: a) whether the content of the dogma can be grasped
in any objective conceptual form whatever, and b) whether the
reality of the dogmatic representation can be tested according to any
of the objective laws known to us.

In both respects, however, it is easy to prove that
the relation of scientific cognition to the dogma cannot possibly
be a positive one, but must be a negative relation;
for a) the concepts ``omnipresence'' and ``limitation of individuality''
contradict each other: an omnipresent
individuality is unthinkable, and everything that
% --- p. 35 ---
\setcounter{footnote}{0}%
\opage{35}can be said about an individual multipresence is---nonsense.
Even if the content of the dogma were not in such evident conflict
with the concept as it really is, it still could
not be appropriated in scientific cognition; for b) the representation
of a glorified individuality that ``fills all
in all'' does not belong in the order of things whose laws
we know, but seeks its object in the supernatural.
But as soon as science engages with the supernatural,
it dissolves itself, for here there is no limit to
fantasy: every profound ``genius'' who appeals
to ``religious experience of life,'' every speculative charlatan
who invokes ``immediate and living intuition of ideas,''
can then without further ado assert whatever he pleases,
and pass off his assertions as scientific cognitions.

It is otherwise with religious cognition. In the religious
sense the word \textit{Credo} in the proposition \textit{Credo, ut intelligam}
takes on a decidedly subjective meaning, and therefore the word
``\textit{intelligam}'' must be taken in the very same meaning. ``I believe in order
that I may know,'' according to Christ's promise (John 8, 31--32):
``If ye continue in my word, then are ye my disciples
indeed; and ye shall know the truth, and the truth
shall make you free.'' This cognition, which does not objectively
reason about the dogma, but in inward and subjective conviction
appropriates the objective dogma, can alone place
itself in a positive relation to the dogma. But it must be
firmly held and not forgotten that, in placing itself, as subjective cognition,
in its positive relation to the definite dogma,
it at the same time places itself in a negative relation to all
possible objective scientific cognition of truth, by whatever
name it may be called.

% --- p. 36 ---
\setcounter{footnote}{0}%
\opage{36}This negative relation is now partly a relation of
indifference, insofar as the believer practically holds to the dogma
without troubling himself about the objective concepts of science;
partly a polemical relation, insofar as the believer also
becomes aware of the objections by which science
may attempt to weaken the religious significance of the dogma. In order
to know the dogma of Christ's seat at the right hand of the Father,
the knower may in the name of faith by no means undertake
to fathom whether the omnipresence of a glorified individuality
can be \emph{thought} or not \emph{thought}. The cognition of faith
remains at \textit{Ita}, and does not advance to \textit{Quare};
it remains in Christ's word, and does not engage in any
speculation.

When, then, after this brief preparation I attempt
what Dr. M. demands, to ``let the theory of cognition go hand
in hand with an examination of the development of the dogmas itself,'' the
question is: what theory of cognition can one then find in Dr.
M.'s development of the dogma of Christ's seat at the right hand of the
Father?

That Dr. M. in the doctrine of ubiquity wishes to correct the old
Lutheran dogmaticians, and that insofar ``the reformatory
moment'' shows itself here too in his Dogmatics, I see
well enough; but I also see what sort of corrective he
has applied. One can namely, on the point treated here,
correct the old ones in two ways, according as one either
appeals to the cognition of faith or to objective
scientific thinking. In the first respect one shows that
the dogmaticians who fixed the notion of the omnipresence of Christ's body
were just on the point of overstepping
the limit of religious cognition, and were saved from objective speculation only insofar as
% --- p. 37 ---
\setcounter{footnote}{0}%
\opage{37}they, by virtue of faith,
let the understanding draw a consequence in which
the understanding itself perished. In the latter respect one concedes,
on the contrary, that the old ones were on the right track in
making ``Christ's seat at the right hand of the Father'' and his individual
presence in the holy Supper an object of
objective thinking, but then seeks to correct their objective
comprehending by a still deeper-going comprehending.
Which way Dr. M. has gone with his ``reformatory'' cannot
be called into doubt: he has evidently gone the latter.

That Dr. M., in ``further'' continuing with his objective
``cognition of the truths of revelation,'' grounded though it be ``in religious experience of life''(!),
further wishes most definitely to have the objective cognition in
faith and out of faith separated from another objective cognition
that is not of faith, I have never overlooked either.
There are thus, according to Dr. M.'s ``fundamental view,'' to be two
kinds of objective cognition, namely: a lower, worldly, general-scientific one,
which cannot yet say \textit{Credo, ut intelligam},
since it is still in its sinfulness, and in this its
sinfulness is only able to give birth to the ``cosmic,'' ``pantheistic,''
``rationalistic'' knowledge for which faith is a paradox; and a higher, ecclesiastical, dogmatic-scientific cognition,
which can at once say \textit{Credo, ut intelligam}, since
it is redeemed by God's grace, and redeemed in such a way that
it can produce the genuine theological speculation for
which faith is no paradox (that is to say, when one
has in addition so much knowledge of languages that one can seek for the word
``paradox'' the corresponding expression in the mother tongue).

How Dr. Martensen wishes this separation to be
% --- p. 38 ---
\setcounter{footnote}{0}%
\opage{38}understood he states most definitely in his ``Dogmatic
Elucidations,'' where among other things (p. 22) he ``submits''
to me ``to consider whether I trust myself to establish
that the redemption of fallen human nature does not
also comprise the redemption of thought, that it is not the purpose of Christianity
to restore true humanity completely,
but to let essential elements---and thought
surely belongs to the essential elements---remain outside
its sanctifying and perfecting activity.''

In the speculative pulpit style, which preaches speculation
without thinking the thought, such an assurance about
``the redemption of thought'' may no doubt look quite well, both
edifying and speculative; but---``show me thy faith
by thy works!''

Were the idea of ``the redemption of thought,'' as
Dr. M. conceives it, to have any true speculative-dogmatic
validity, then this redemption would surely have to be for the dogmatician
what genius is for the philosopher. It would
then have to turn out that the dogmatician whose ``genius'' had been
``redeemed'' was really able to treat the problems of revelation
with such scientific profundity, and
to develop the concepts with such acuity, that even the
most gifted philosopher would have to envy ``the redeemed one''
not only his ``capacity for immediate and living intuition of ideas,''
but also his wondrous gift for higher thinking
and penetrating comprehension. Is this now in the remotest
degree applicable to the Martensenian treatment of the
Christian dogmas? I submit to the reader to consider
whether he ``trusts himself to establish that Dr. M.'s scientific
development of the dogma of ubiquity bears the least trace

% --- p. 39 ---
\setcounter{footnote}{0}%
\opage{39}of any supernatural ``redemption of thought.'' For I really
cannot find that this ``dogmatic genius,'' either in §~177
or in any other dogmatic paragraph, has made the slightest
scientific provision for rendering the invoked separation of
the ``presence of Christ'' and the ``presence of the Logos''
in any degree evident to objective thinking. On the contrary,
it seems to me that everything he ``produces'' on this
occasion is of such a nature that not even the simplest
``dilettante in logic'' will envy him the ``productivity.''


``But'' --- the ``well-disposed'' reader perhaps objects ---
``what you say here may, from your standpoint, be quite
correct; but what does it prove? It proves, after all, only
what Dr. Martensen too (Dogm. Oplysn. p. 45) can testify:
that from the beginning you have misunderstood \textit{Credo},
\textit{ut intelligam}. Since you, then, as Dr. Martensen can
testify, so frivolously misunderstood \textit{Credo, ut intelligam}
from the beginning that, without thinking about your faith
--- \textit{quod negligentiæ mihi videtur}! --- you even sat and
waited for ``an inheritance of ontological and other cognitions
to fall to you by a stroke of luck'': what wonder, then,
that in your ``rashness'' and --- I must tell you this seriously
--- in your ``frivolity'' you cannot see what is significant
in Dr. Martensen's productivity: to cognize what redeemed
thinking produces here, your own thought would of course
itself have to be redeemed. We others, on the other hand, who
from childhood have been brought up, and precisely thereby
reborn, in Christianity, we others who grew used to the paradox
already when we were weaned from our mother's breast, all we
others who have thus little by little become familiar with the
% --- p. 40 ---
\setcounter{footnote}{0}%
\opage{40}truths of revelation, as we have thus little by little received
the ``redemption of thought,'' see, all we others, we can testify
for Dr. Martensen that he, at least relatively and to a certain
degree, has comprehended both the presence of the Logos and the
presence of Christ, both the presence of Christ at the right
hand of the Father and the presence of Christ in the Holy Supper;
indeed, all we others can for our part testify that Dr. Martensen
has speculatively understood Christianity, just as Dr. Martensen
will surely also for his part be able to testify that all we
others have speculatively understood him!''

Dear reader, whoever you may be, reborn or not reborn
(for about that I can testify nothing in the present case),
before you grow any more enthusiastic about speculation's
\emph{testifying} method, permit me just one small question:
Do you really trust yourself to demonstrate that Dr. M., in
any dogmatic paragraph whatsoever, and particularly (to dwell
on this single dogma) in the paragraph on Christ's omnipresence,
has had the courage to take on the consequence of his own
speculative testimony?

Had Dr. M. in earnest himself believed in the ``redemption''
of objective scientific thought, he would of course have
had to treat the ``presence of Christ'' in quite another manner
than has happened. He would then not merely have had to
``testify'' that the concept of the presence of Christ in
heaven and of the individual presence of Christ in the Supper
is ``to the natural man a paradox,'' but ought also to have
``testified'' expressly that he for his part could very well
--- at least to a certain degree --- think this supernatural
concept speculatively by the help ``of the redemption of thought'';
indeed, he ought to have testified that both Anselm (the man
who early
% --- p. 41 ---
\setcounter{footnote}{0}%
\opage{41}``inspired in him the desire for dogmatic speculation'') and
``the first Adam'' and ``Peter on the Mount of Transfiguration''
really have been able to think what is otherwise unthinkable to
the human understanding, the concept of the individual
multipresence of Christ's body, and --- mark well --- to think
this concept scientifically, objectively, contemplatively,
speculatively, dogmatically. In short: Dr. M.'s appeal to the
``redemption of thought'' contains the claim that Christ and
the apostles are supposed to have promised the faithful a new
speculative logic, an objective logic or ontology, whereby later
on, when the first evangelical miracles had ceased, it might become possible for a ``dogmatic genius,'' as ``the organ of
the holy and catholic Church,'' to perform a new and more
up-to-date miracle, a speculative scientific miracle, namely by
making that thinkable which every other scientific thinker must
otherwise find unthinkable.

Now has Dr. M., when it came to the point, really been
in earnest with this claim? I cannot see it. For the little
thought-experiment which the ``dogmatic genius'' conducts here
on the omnipresence of Christ's body does not exactly seem to
me to display any ``sanctifying and perfecting activity'' or
to bear the trace of any peculiar logic of rebirth.
--- ``Even a transfigured individuality, even a spiritual body,
cannot be thought without limitation.'' This thought is really
so modest, so accessible, so straightforward, natural and meek,
that it does not in the least exceed what ordinary logic can
manage in this respect.

When, now, Dr. M. with all this complains so bitterly that his reviewers have not ``taken the time to let
% --- p. 42 ---
\setcounter{footnote}{0}%
\opage{42}the assessment of the doctrine of cognition go hand in hand with
an examination of the development of dogma itself,'' and wants
the reader to regard the introduction's ``general principles in
indissoluble connection with the actual view of the matter,''
claiming that these general principles first acquire ``life and
shape'' through such a view of the matter, then I really do not
see what advantage the Right Reverend author can expect from
insisting on such a demand. For, I will not conceal it, it seems
to me as if Dr. M.'s ``doctrine of cognition'' looks far better
when it is kept at a certain general and hovering distance from
the ``actual view of the matter'' than when one applies it to
``an examination of the development of dogma itself.'' What His Rev.,
e.g., elucidates in the Dogmatic Elucidations about ``the immediate
and living intuition of ideas,'' about the ``redemption of thought,''
about the ``consciousness of redemption'' and the ``consciousness
of revelation,'' about all ``the glory and divinity and heavenly
clarity'' which ``revelation'' by means of the ``consciousness of
revelation'' has in itself ``on account of the insight it lets
us gain into the depths of God,'' which insight is in turn
confirmed by speculative glances toward Anselm and the first Adam
and Peter on the Mount of Transfiguration, and this in such a way
that these glances are in turn confirmed by new insights into
``the prophetic vision of the new heaven and the new earth,'' into
``the transfiguration of corporeality,'' into ``the rose of
contemplation'' and the ``quiet joy'' of speculation (``In thy light
shall we see''), all this seems to me to be set forth so beautifully,
so profoundly, wittily and upliftingly, that it is almost a sin
to drag it down from the ethereal height of the hovering and put
it in connection with any ``actual
% --- p. 43 ---
\setcounter{footnote}{0}%
\opage{43}view of the matter''; for through the connection with something
as prosaic as an actual view of the matter all these ``roses of
contemplation'' would of course fade and lose their fragrance.

How, then, it will fare with the Martensenian ``doctrine
of cognition'' when it really has to ``go hand in hand with
an examination of the development of dogma itself'' has now
been shown in the development of the dogma of ubiquity. So that
no one may think that I have here arbitrarily chosen a dogma
which by its nature belongs among the most incomprehensible,
I must point out that I in no way acknowledge the recently
refreshed distinction between ``more and less comprehensible''
dogmas, and that Dr. M., whose acuteness deserves the credit
for this distinction, must himself concede that this dogma,
according to his own ranking, belongs among the most comprehensible.
``The incomprehensibility of darkness'' --- he says, Dogm. Opl.
p. 58 --- ``is essentially different from that of light, from
which it also follows that the comprehensibility of which there
can be talk here is a specifically different one\dots{}''; but the
transfigured Christ belongs, after all, not to the kingdom of
darkness but to that of light. Dr. M. therefore cannot complain
that on this occasion I confuse the ``teleological'' with the
``anti-teleological'' for him; for ``Christ's royal office'' surely
belongs to ``the teleological.'' --- What, however, has chiefly
led me to dwell so long on the dogma of ubiquity is by no means
the Martensenian ``opposition between the teleological and the
anti-teleological,'' but, as I said, the Stillingian note.

In thus concluding the discussion of Mr. Stilling's
competence as reviewer of Dr. Martensen's Dogmatics, I shall
permit myself only one more remark.

% --- p. 44 ---
\setcounter{footnote}{0}%
\opage{44}That Dr. M., if he felt disposed at all to answer ``his
reviewers,'' would have to answer in a sharp tone was to be
expected, given ``the language'' one had permitted oneself toward
His Rev. That Dr. M., at the moment he condescends to enter
the arena, at once shows his opponents that there are birch
twigs in the broom with which he now intends ``to sweep before
his own door'' is polemically in order.\footnote[1]{Thus it is, e.g., quite fitting and by no means unwitty that in his Dogmatic Elucidations he presents me as the one who wants to inherit from the ``rich uncles.'' ---} That Dr. M. does not stop
at the intellectual birch, but (just as if he involuntarily felt
that his intellectual birch was not salted enough) also seizes
the moral (sit \textit{venia verbo}) broomstick, in order to act
with all the greater ``earnestness,'' is perhaps not quite in
order; still, since Dr. M. has once thought it his duty to
execute the moral chastisement, this prerogative too shall, for
my part, be willingly granted to his ``dogmatic genius,''
and all that he has said on this occasion about my ``obtrusiveness,''
``rashness,'' ``frivolity'' and ``inconstancy'' forgiven for the
sake of old friendship and acquaintance.

On the other hand there is one thing which must most
seriously be held up to Dr. Martensen, and that not merely for Dr. Martensen's own sake but also for the sake of ``literature''
(for Dr. Martensen and ``literature'' and ``literature'' and
Dr. Martensen cannot be separated here), and it is that
His Rev. should not another time, when he wishes to discuss
``the general principles,'' be in a hurry
to \emph{lay} his opponent's ``scribbling'' aside
% --- p. 45 ---
\setcounter{footnote}{0}%
% CHECK p. 45: centred line: heading or verse
\opage{45}\begin{center}15\end{center}


before he has properly looked at what it was that stood
written in this ``scribbling.'' That Dr. M. on this
occasion thought he ought to ``keep justice in literature,''
I do not doubt; but --- for another time's sake ---
it would really be desirable if a man in whose authority
Dr. Martensen places confidence would properly impress upon him
the above well-meant advice; for if His Rev. persists with
his excommunicating language of power, as he has begun here,
he might easily misuse the thunderbolt of the ban to great harm
both for himself and for ``literature.''

As a kind of excuse for the dismissal of Mr. Stilling
one might perhaps regard Dr. M.'s remark
(Dogm. Opl. p. 8) that he ``cannot possibly take account of
everyone, already for the reason that this would lead to a tiresome
and unmanageable prolixity.'' On the assumption
that everything is settled when the ``propositions'' on which my
``examining review'' rests are overthrown, the Right Reverend
author thinks, then, that he hits the nail on the head
by hitting my head, and so it is now I alone
who must bear the brunt.

Without wishing to dwell on my ``manner of writing'' (which
may also well be regarded as superfluous, after His Rev. has himself
given ``proof of his manner of writing''), Dr. M. wishes merely to
dwell a little on the ``form and tone'' of my review,
and then goes straight at the matter, declaring:

``The propositions which the reviewer (Prof. N.) sets up against
me I regard as propositions for which no one but
himself is responsible. Whether those propositions perhaps have
an entirely different meaning in the humorous context
% --- p. 46 ---
\setcounter{footnote}{0}%
\opage{46}in the pseudonyms, whether the reviewer here is in understanding,
or in non-understanding, or in misunderstanding, all this lies
outside the circumference of my investigation. I shall here only
confine myself to the far lesser and more modest task
of sweeping before my own door by showing that I have no
use whatever for the doctrines which the reviewer with so much
obtrusiveness has wished to impart to me, since I must
apply to these the old saying that what is true in them is not
new, and what is new not true; just as I shall also show that,
however it may stand with his understanding
of the pseudonyms, it is irrefutably certain that he has not
understood me, but has most arbitrarily misunderstood and
misinterpreted my thoughts'' (Dogm. Oplysn. p. 13 ---14).

That my ``examining review'' is at certain points
somewhat pointed, I will not deny; but the ``well-disposed''
reader will then surely not deny either that the review's merely
juxtaposing, examining and questioning form is lenient
and weak in relation to the strong words in which
Dr. Martensen here proclaims the unassailability of his dogmatic
speculation.

The review, which begins with ``an apology for permitting
itself to juxtapose two writings so heterogeneous in direction,
spirit and tone,'' as Mr. S. Kierkegaard's
``Johannes Climacus'' and Dr. H. Martensen's ``Christian
Dogmatics,'' divides into two parts.

The first part of the review (pp. 9—12) contains only
what Dr. M. calls ``an excerpt'' from the pseudonym, a condensed
presentation of the main thoughts that move the humorist,
as he takes the liberty of treating the theologians'
attempts at objective (historical and speculative)
% CHECK p. 46: note mark without a note

% --- p. 47 ---
\setcounter{footnote}{0}%
\opage{47}knowledge of the objective truth of Christianity, not as a
scientific task that properly ought to be discussed in science,
but as the sort of ``earnest'' endeavors which, when
they are seen in the idea, prove to be of such a nature that
they essentially belong in the higher comedy. As for the
comic execution itself, it naturally could not
be reproduced in an ``excerpt''; which is why I have not permitted
myself any attempt in this direction either. But that a
terrible earnestness runs through the ``humorist's'' comedy was
at the same time so clear to me that I, especially when I saw how

Dr. Martensen in the ``Preface'' to his Dogmatics had hit
upon, and spoke so significantly of, ``stray thoughts
and aphorisms, fancies and flashes,'' could not possibly
refrain from directing His Rev.'s attention to the
question: \emph{Has not ``Johannes Climacus'' dialectically
sharpened the Christian problem?}

Since, then, the first part of the review occupies itself
exclusively with the pseudonym's theses, and does not with a
single word enter into the Martensenian train of thought, the
proof that I have misunderstood Dr. M. cannot be fetched from
this section. The whole emphasis therefore falls on the review's
second part, which (pp. 42 --- 132) throughout turns
on the question:

\emph{Has not ``the Christian Dogmatics'' undialectically
evaded the Christian problem?}

In this second and comparatively more detailed division (which
is even separated from the first by a little rule) I have,
as far as I know, not used any foreign
tongue either, but, with the exception of quotations which do not
readily admit
% --- p. 48 ---
\setcounter{footnote}{0}%
\opage{48}of any rewording, set forth everything in my own words, and
explained to Dr. M. my opinion of his Dogmatics as well as
I have been able.

If Dr. M., then, wants to make it ``irrefutably certain'' that I
have not understood him, but ``have most arbitrarily misunderstood
and misinterpreted his thoughts,'' then I am surely also entitled
to expect that His Rev. will partly enter into
the main points developed in my review, partly ground
his own ``misinterpreted'' concept of science so clearly and definitely
that it can be seen from this that my conception has been a misunderstanding.
Perhaps Dr. M. will one day in time do both;
but that he has not yet in his ``Dogmatic Elucidations''
done either the first or the last is --- irrefutably
certain.
\emph{Dr. Martensen has, after all, in his Dogmatic Elucidations not at all undertaken to answer}
the objections which my examining review
puts forward against his Dogmatics.

In order at once to indicate the difference between the cognition of faith
and objective knowledge, the second part of the review
begins with an investigation of what it means to be
``absolute in faith'' and ``absolute in knowledge.'' The author of the
dogm. Opl., however, is so far from taking any account
of this investigation that he even ascribes to me such
opinions as he could not possibly have ascribed to me
had he in the least heeded my words.

As regards knowledge, e.g., it is expressly said
(Anm. p. 49): ``To declare oneself absolute in knowledge is a presumption,
insofar as the one who speculates means thereby to claim to
have a knowledge, absolute and perfect in all respects, of
% --- p. 49 ---
\setcounter{footnote}{0}%
\opage{49}the absolute\dots{}'' Nevertheless Dr. M. acts as if
such a thing had never been said, and imputes to me without further ado the
presumption of wishing to proceed from a speculation that undertakes
to make God's essence absolutely comprehensible. --- ``But,
answers Prof. N., .... God must either be absolutely comprehensible
or absolutely incomprehensible.'' (Dogm. Opl. p. 52).

As regards faith, the review says (p. 50): ``To
wish to pass for absolute in faith is a presumption
if the individual means thereby to give himself out as able to fulfill
all the work of faith, to have a faith absolute and unconditional
in all ways, the faith that moves mountains.''
But, far from entering upon the distinction emphasized in the review
between holding absolutely to the principle of faith
and giving oneself out as an absolute believer (a distinction
which the preface to the Dogmatics made highly necessary),
the author of dogm. Opl. passes this point by with apparent
ignorance, and ascribes to me without further ado the presumption
of wishing to give myself out as a hero of faith (cf. dogm. Opl. pp.
71—72 and elsewhere), and this His Rev. has, as it seems,
done solely and alone in order to be able (Opl. p. 12) to have the
pleasure of repeating the ambiguity he has placed in
the preface to his Dogmatics, ``that just as there was a time when
we had not a few among us who had become absolute
in knowledge all too early, so there is now no lack of such
among us as have become absolute in faith all too early\footnote[1]{That a Mr. Cand. Hjort can come running with the news that I now really imagine myself to be the believer himself, because in my writing ``The Faith of the Gospel and the Modern Consciousness'' I have let the believer speak, I must naturally put up with; but that Dr. Martensen now also takes this discovery to his credit weighs a little heavily on me. Dr. M. ought surely to be able to distinguish between letting the believer speak and speaking as the believer; Dr. M. himself says, moreover, that he has read that writing of mine whose preface ends thus: ``The author of the present writing sincerely confesses that he does not yet, either in word or in deed, rightly understand how to contend for the faith as is fitting'' etc.} % the note runs over onto p. 50
and have become defenders of the faith all too early.''
% --- p. 50 ---
\setcounter{footnote}{0}%
\opage{50}As a contribution to a closer determination of the relation
between the cognition of faith and science, the review treats, under the heading:

``The Christian Dogmatics'' gropes for the problem, but

% CHECK p. 50: centred line: heading or verse
\begin{center}does not find it,\end{center}

the question whether speculative dogmatics can essentially
be distinguished from so-called philosophy of religion
or not. The answer to this question is of thoroughgoing
importance in the present matter. If it really could
be shown that, in respect of principle, method and scientific
interest, there existed a true and significant opposition
between the speculative dogmatician's and the speculative
philosopher's conception and treatment of the Christian
dogmas, then Dr. Martensen would in truth be right; then
one would really, after all, finally be compelled to admit that there
was a double speculation, a dogmatic and a philosophical,
and that the first was in its kind just as scientifically
justified as the last. Now the ``review'' strives to demonstrate
that such an opposition exists only in the dogmatician's
imagination, not in the matter itself. The ``review'' shows
that the attempt Dr. M. has made at separating the concept of dogmatics
from that of philosophy of religion dissolves --- despite
all his assurances and claims --- into a pure nothing;
indeed, it is expressly emphasized that the Martensenian separation of
% --- p. 51 ---
\setcounter{footnote}{0}%
\opage{51}``theological'' and ``philosophical speculation'' has had so
unfortunate an outcome that even two paragraphs in his
own Dogmatics have on this subject fallen into open
enmity with each other.

What, now, has the author of dogm. Oplys. elucidated on this
occasion? Nothing, absolutely nothing. Without heeding the objections
put forward against his separation, he merely repeats
his claim that theological speculation is of an entirely
different and higher kind than the philosophical, forces upon me, or
rather fabricates for me (Opl. p. 14), the assumption that by
objective knowledge one should think merely of ``an indifferent
knowledge,'' although the ``review''\footnote[1]{``But when it is now well known that this separation (between the dogmatician's and the philosopher's interest in the cognition of dogma) rests on an unhappy confusion of the personal and that which pertains to genius; when it is now well known that one person can be a believing Christian and work vigorously in faith without having the intellectual, i.e. the theoretical, love of which there is talk here, while another can be wholly governed and captivated by this love without to the same degree working for the subjective appropriation in faith; when it is well known\dots{}. that the intellectual and speculative love with which the thinker immerses himself in the mysteries of revelation is by no means any properly religious requirement (``faith is, after all, the one thing needful for us all''), but solely and alone the urge and need of talent for the contemplation of the idea, for objective joy in divine visions: where then is the difference?\hfill Unders. Anm. p. 65.} says precisely the opposite; % print: at,En (misprint)
indeed, he even (Oplys. p. 27) takes me sternly to task,
because I am supposed to have overlooked that every knowledge must
have its interest, since ``an absolutely disinterested knowledge can hardly
be thought.'' --- Thus, now, Dr. M. understands his
``reviewers.'' One draws his attention to the difference
between the religious and the scientific interest; from this
he infers that one is speaking of ``an absolutely disinterested

% --- p. 52 ---
\setcounter{footnote}{0}%
\opage{52}knowledge,'' and then complains that one ``misunderstands and misinterprets
his thoughts.''

``The Christian Dogmatics wishes to appropriate speculation,
but without grasping the problem of speculation.''

Under this heading the ``review'' (pp. 67—90) enters
upon a closer development of the principal difference
between the cognition of faith and speculation which was pointed out in the foregoing. The investigation does not here apply any
foreign standard, but everywhere finds its point of departure in
Dr. M.'s own Dogmatics; it is Dr. M.'s own ``principles''
concerning cognition in faith and out of faith that, with
a careful juxtaposition of the relevant quotations,
are made the object of scientific examination. After
having secured the concept of what Dr. M. calls ``pure
science,'' the ``review'' concentrates its whole attention
on the passage in the ``Dogmatics'' which, as regards the standpoint,
may rightly be designated the real main passage.


``For dogmatics the absolute truth of Christianity is
given in advance and independently of all speculation. That
δὸς ποῦ στῶ which the seeking philosophy has so often
uttered is originally answered for dogmatics; and the
dogmatician does not make Christianity dependent on his inquiry,
but only seeks through thinking a deeper appropriation
of the truth which is to him the most certain of all, and to
which he has come by an entirely different way than that of speculation.''
(Dogm. p. 5.)

The hidden contradiction which lies in these words the
review draws out, and concentrates its result in the following expression
(p. 80):\dots{} ``is the thinking of faith different from
% --- p. 53 ---
\setcounter{footnote}{0}%
\opage{53}speculation's, or are they at bottom one? If the dogmatician grants
the latter, then we admit that he is right
when he seeks the completion of the cognition of faith in speculation;
but then we must ask what entitles him
to the claim that he himself has come to the highest
truth by an entirely different way than that of speculation? If, on the other hand,
he declares for the first assumption, namely that
the cognition of faith is in principle and essence entirely different
from speculative cognition, then we gladly grant him
that he must have come in advance to the absolute truth
by an entirely different way than that of speculation; but then we must
ask where, as dogmatician, he then wants to go with speculation?''


What, now, has the author of dogm. Opl. done to
elucidate this point? He has (Opl. p. 51) foisted on me
the following blunder:

``As regards, then, Prof. N.'s claim that I
have misunderstood the problem of speculation, this evidently
rests on his setting up an entirely different concept of
speculation than the one I have acknowledged in my Dogmatics,
and thereby from the very first displacing the whole point of view.''

Thus: I keep to Dr. M.'s own concept of speculation ``set up''
in the Dogmatics, in order to show that this concept of speculation
is in every respect false and untenable; without
entering into my reasons, Dr. M. declares that I
set up ``an entirely different concept of speculation than the one
he has acknowledged,'' and then wants this declaration
to make it ``irrefutably certain'' that I ``have not understood him,
but have most arbitrarily misinterpreted and misunderstood his
thoughts.''

% --- p. 54 ---
\setcounter{footnote}{0}%
\opage{54}Indeed, the author of dogm. Opl. goes a little further still.
He continues, namely, to raise accusations which are of such a
nature that I must almost assume that His Rev. has either
read my review as he, by his own admission, has
read the pseudonyms, i.e. merely cursorily and ``fragmentarily,''
or else forgotten again what he had read.

On the occasion that I (which, as shown above, is
an untruthful claim) am supposed to have proceeded from a speculation
which undertakes to make ``divine things
absolutely comprehensible,'' Dr. M. thus ascribes to me (Opl.
p. 52) the following argument: ``God must either be
absolutely comprehensible or absolutely incomprehensible. For Kant and
Hegel were systematic heads, and the one of these taught
that divine things were absolutely incomprehensible, the
other that they were absolutely comprehensible. If we therefore wish to be
systematic heads, we must teach either the one or
the other.''

What I said in my review (p. 124) is,
however, only what has long been recognized in science:
that he who wishes to write a system and pass for a systematic
head must hold fast to one principle, and from first to last
respect its consequences, whether he is otherwise
a theologian or a philosopher. To indicate what I mean
by the expression ``a systematic head,'' I cited Kant and Hegel
by way of example. Dr. M. takes occasion from this
to reprimand me for a \textit{metabasis in aliud genus}.

``But'' --- it says (Opl. p. 52) --- ``instead of
appealing to Kant and Hegel, whose invocation in the present
context is a \textit{metabasis in aliud genus},
% --- p. 55 ---
\setcounter{footnote}{0}%
\opage{55}he (Prof. N.) should rather have brought forward to me the example of dogmaticians
of corresponding authority.'' ---

The reader who has the slightest acquaintance with the theological
literature of recent times will no doubt be taken aback and fall
to wondering at the peculiar boldness with which Dr.
M. here speaks of \textit{metabasis in aliud genus}. It is, after all,
surely well known, indeed --- ``irrefutably certain'' --- that the dogmaticians
Daub and Marheineke essentially kept to the same
\textit{genus cognoscendi} as the philosophers Schelling and Hegel,
and that the dogmatician Schleiermacher was in his way just as
systematic-subjective as the philosopher Kant.

It must furthermore arouse wonder that the author of
dogm. Opl. can thus blame me as one who has not
``brought forward to him the example of dogmaticians of corresponding authority,''
since I did, after all, in my review (p. 89) expressly
cite Daub and Marheineke to him, and, p. 66, reminded him
of Schleiermacher.

No less striking seems to me the suffisance
with which Dr. M., without in the remotest respect heeding
what stands in the review (pp. 82---84), puts forward the
demand that I ``should in deed have refuted the
proposition he has set up in his Dogmatics, that just as those
have always been put to shame who have given themselves out as
having comprehended everything, so those have been no less put
to shame who have thought once and for all to be able to stake out
the limits of all human comprehending; for afterwards
it always turned out that there was a \textit{plus ultra}, and the supposedly
fixed limits revealed their movable nature
by shifting'' (Dogm. § 33).

``Yes'' --- the review answers (p. 82) --- ``by
% --- p. 56 ---
\setcounter{footnote}{0}%
\opage{56}shifting!'' and then adds a good deal about the ``fluid limit''
and about ``\textit{plus ultra}.''

If ``finding oneself'' and remaining true to one's standpoint
is to show itself in this, that, whatever may be objected
against the ``fundamental view'' once set up, one merely keeps
a stiff upper lip and defies all objections by a bold
repetition of one's own claims, then I gladly admit that
Dr. M.'s conduct in this matter can be called dogmatically normal,
and that His Rev. in this respect must be said to have
``found himself.'' If, on the other hand, one is to lay the least weight
on truth and the investigation of truth, then I really do not see
what Dr. M. can think to achieve by so confident
a repetition of claims whose untenability is so
conspicuous that they find no support even in popular
opinion.

That those who have thought once and for all to be able to
``stake out the limits of all human comprehending'' have always
``been put to shame'' is a proposition as indisputable as it is trivial,
when the talk is of such things as really
can be the object of ``human comprehending.'' Natural science,
medicine, mathematics, philosophy, % print: Mathemathiken (misprint)
in short, all actual sciences are, after all, in a constant
advance; in all actual sciences there is a \textit{plus}
\textit{ultra}, and he would have to be deranged who here would undertake
once and for all to ``stake out the limit of all human
comprehending.''

But who does not now see that in this matter there is talk of
``staking out'' an entirely different ``limit.'' Here there is talk, after all, of
``staking out'' the limit by which imagined knowledge
necessarily distinguishes itself from true knowledge. This
% --- p. 57 ---
\setcounter{footnote}{0}%
\opage{57}limit is very easy to stake out: Of the essence of the supernatural
there is only an imagined knowledge; of the law according to which
the supernatural acts upon the natural there is
only an imagined knowledge; of the objective validity of miracle
there is only an imagined knowledge; of the objective truth of revealed
religion there is only an imagined knowledge.

That the ``limit staked out'' here stands fast is demonstrated
as much by the nature of the matter as by the testimony of history.
It lies in the essence of science that it cannot venture
into the precincts of the supernatural without surrendering itself
to fantasy and fantasticality. It is seen from the history
of theology that this science, far from advancing
in objective metaphysical knowledge, rather goes backward in
the same proportion as the actual sciences go forward.
Natural science, philosophy and historical criticism have,
after all, precisely in recent times, each in its own way, brought theology
into such a self-contradiction that no great acuteness
is exactly needed to see that the theologians' medieval\footnote[1]{In this respect it is, then, also not without significance that Dr. M., in order to find a kind of ``prototype'' for his ``dogmatic speculation,'' goes back to Anselm and the monasteries of the Middle Ages.}
concept of science is dissolving day by day, and shows
itself equally useless for science and for --- the Church.

``The Christian Dogmatics wishes to appropriate faith, but
without holding fast to the problem of faith.''

To justify this claim, the ``review'' (pp. 90—125)
keeps so carefully and exactly to
the ``principles'' which the Martensenian Dogmatics ``sets up''
concerning faith that from the whole dogmatic work one shall not
% --- p. 58 ---
\setcounter{footnote}{0}%
\opage{58}be able to cite to me a single passage whose significance on this
point has been overlooked. It is altogether quite characteristic
that the Martensenian Dogmatics, whose ``speculation'' is, after all, in
faith, proceeds from faith and aims at faith, betrays precisely
what Mr. Stilling in his review (p. 29)
aptly remarks --- ``a striking lack of qualitatively
designating predicates for faith and hope.'' This,
I say, is characteristic; for had Dr. M. really become
attentive to that which properly constitutes the quality of faith,
he would necessarily have had to give up ``speculation.''

That faith unites in itself both the subjective and the
objective, that the believer must keep to the object of faith,
must know what it is he believes, and give himself an account
of the content of his faith, is self-evident, and has not
been disputed either. But precisely as both sides
grant that faith and the object of faith must correspond to
each other, there recurs here the question which in
the present matter is the real main question: How
is the relation between faith and the object of faith
to be rightly determined?

In Dr. M.'s Dogmatics the relation is determined at
once as a relative relation of knowledge and, besides, at the same time as
an existential relation, i.e. a subjective personal relation.

That this conception of the relation is unclear, ambiguous and
self-contradictory, the review strives to show, not by applying
a foreign and arbitrary standard, but by keeping % print: vilkaarl= (misprint)
to the Dogmatics' own words.

``The second main proposition of dogmatics'' --- it says, Anm.
p. 96 --- ``aims at preventing `the misdirection which,
out of concern for `the believer' and his state of soul, becomes
% --- p. 59 ---
\setcounter{footnote}{0}%
\opage{59}indifferent to faith (fides, \textit{quæ creditur}), which out of zeal for
the edifying forgets the content with which one is to build,
forgets the foundation on which one is to build.'\,'' --- ``What
the Reformation wanted'' --- it says (Dogm. p. 40 § 21) ---
``was neither exclusively the objective nor the subjective;
it was the free union of the objective and the subjective,
of the content of faith and the inwardness of faith, of divine
revelation and religious self-consciousness.''

Against this ``free union of the subjective and the objective,''
as Dr. M. has understood it, the review raises
the double objection, namely that a) ``the objective'' thereby
dissolves, so that what remains is only a vague and
indeterminate subjective something, and b) that ``the subjective'' on its
side is likewise volatilized in such a way that what remains
is only an indeterminate and indeterminable objective something.

a) As regards the first point, Dr. M. naturally wants
to secure ``the objective'' by a scientific appeal
to ``the testimony of the Church, of Scripture and of the Spirit.''
The review by no means reproaches Dr. M. for having
referred us in his Dogmatics to these testimonies; but
it draws his attention to the fact that he has not understood how to
make use of these in such a way that either faith itself (\textit{fides, qua creditur}) or the content of faith (the doctrinal concept --- \textit{fides, quæ}
\textit{creditur}) is thereby in the least clarified and secured. The review's
train of thought is in brief the following:

``The objective canon for all Christianity'' (Dogm. §
22) is Christ himself. But Christ is only in his Church:
therefore it is the Church, ``that great individual!'' --- that is the
objective canon. But the Church's real ``touchstone (lapis
lydius)'' is only in Scripture: therefore it is Scripture that is
% --- p. 60 ---
\setcounter{footnote}{0}%
\opage{60}the objective canon. But the right understanding of Scripture is
only in the testimony of the Spirit: therefore it is the testimony of the Spirit
that is the objective --- no, that is not true ---
the testimony of the Spirit is an absolutely subjective and ``inward
canon.''

All that is said here about ``the objective'' thus goes
back into the entirely subjective; for the testimony of the Spirit
is in the subjective, not partly in the subjective and partly
in the objective, but solely and alone in the subjective, absolutely
in the subjective. How easy it is to see this, if
only one will see. All theologians, the rationalists and the
speculative ones as well as the orthodox, are, after all, in the Church,
want, after all, to be in the Church, appeal to the Church and defend
the Church, and that although their concepts of Christ
as well as of Scripture and the Church are as contradictory and conflicting
as can be conceived. If the Church --- ``that great
individual''! --- is now really to be assumed to embrace, acknowledge and appropriate
all these mutually heresy-hunting and condemning
tendencies, then the testimony of the Spirit in ``the holy
and catholic Church'' is certainly not merely a many-sided, but also
so ambiguous and double-tongued a testimony that he who
is rightly to interpret this testimony must, after all, have a still
higher ``testimony of the Spirit.''

``And what'' --- asks Dr. M. (Dogm. p. 57) ---
``was rationalism other than a great theologia
irregenitorum which flooded Protestant Christendom?''
The objective question of the difference between true
and false theology, which in dogmatics was to be decided scientifically,
is thus after all imperceptibly transformed into a subjective
and unscientific question of which theologians are
% --- p. 61 ---
\setcounter{footnote}{0}%
\opage{61}``reborn'' and which are not reborn. See, here we have
the vague, indeterminate subjective something which must nevertheless in the end
turn the scale when one has taken the object of faith in vain by treating it as another straightforwardly scientific
object, and precisely thereby lost all objective marks of distinction.

What, now, has the author of dogm. Oplysn. done to
elucidate this point? Put forward (Oplysn. p. 16) the ingenious
remark that ``even the poorest rationalist dogmatics
must recognize that the truth has a certain objectivity''..,
and besides --- though without citing any passage
that has been passed over --- (Dogm. Oplysn. p. 41) tossed out the
accusation that one ``does not at all enter into the opponent's
train of thought,'' does not investigate ``the genesis of the concept of faith
in the author against whom one directs one's attack,'' but
``without further ado assumes another concept of faith as the true one,
deals with it as a fixed presupposition that is raised
above all doubt'' etc.

b) Concerning the second point, the review by no means overlooks
that Dr. M. also wants to assert ``the subjective,''
that he calls the cognition of faith an existential cognition, and
wants ``the religious relation to God to be regarded as an existential relation'';
but --- \textit{hinc illæ lacrymæ}! precisely because Dr. M.
wants to appeal to existential thinking, precisely therefore
``the critique'' reproaches him for occupying himself with clarifying
the content of this thinking in the objective knowledge of speculation;
for speculative thinking is in principle and essence the exact
opposite of existential thinking. For existential thinking
the subjective, the personal, is the absolutely decisive; in
speculative thinking the thinker goes out of himself,
% --- p. 62 ---
\setcounter{footnote}{0}%
\opage{62}loses himself in his object and relates himself to it as objectively
and impersonally as is at all possible.

``When the dogmatician'' --- says the review, p. 101, ---
``comes onto the subjective track, he recognizes and inculcates
subjectivity in the best way; when he then
falls over again into the speculative train of thought, he asserts
objectivity, objective knowledge, `theological speculation,'
most emphatically; and yet these two directions are
not merely entirely opposed to each other, but also entirely
contradictory.''

Dr. M., who to a certain degree --- but also only relatively
and to a certain degree --- wants to have ``the subjective,'' and
likewise to a certain degree --- but also ``only relatively and
to a certain degree'' --- wants to have the objective, thus gets neither
the objective nor the subjective decision, but a neutralizing
union which is --- unprincipled.

What, now, has the author of dogm. Oplysn. done to
elucidate this point? Has he anywhere attempted to make clear
to us how a cognition which is as subjective,
as inward, as absolutely personal as a true cognition of faith
must necessarily be, can be converted into the objective,
impersonal form of speculation and science, and still
be cognition of faith? No, His Rev. has come far
more easily by his dogmatic elucidation, since he (p. 42)
has pointed out that what I call existential
thinking must presumably be about ``the same as what he too in
his Dogmatics has called existential thinking and,
with Daub and Mynster, called cognition in religion.''

Thus, then, this matter now stands. One has permitted oneself
to review ``the Christian Dogmatics,'' whose preface
% --- p. 63 ---
\setcounter{footnote}{0}%
\opage{63}looked so inviting precisely in this respect. As regards the
two reviews treated here (Mr. St.'s and mine):
Dr. Martensen dismisses the first, does not --- as has been
shown --- enter into the last, but declares outright
that the critique of his Dogmatics ``cannot yet really be
said to have begun.'' One has in these reviews put to
Dr. M. definite theses, shown him definite contradictions
in his Dogmatics, cited definite reasons to him. Dr. M.
declares the whole to be an empty ``formalistry'' and appeals
to the material.

But although Dr. M., despite all his definiteness, as
said, will not enter into any definite answer to
any definite main question, he nevertheless wants
to make it ``irrefutably certain that one has not understood him,
but has most arbitrarily misunderstood and misinterpreted his
thoughts.'' --- The discussion thus takes the following course:

One draws Dr. M.'s attention to the fact that in the cognition of faith
it becomes necessary to hold absolutely to the principle
of faith: His Rev. understands it as if one wanted to give oneself out
as ``absolute in faith.''

One draws Dr. M.'s attention to the fact that in science one
must hold absolutely to the principle of knowledge. His Rev.
understands it as if one wanted oneself to give oneself out as ``absolute
in knowledge.''

One draws Dr. M.'s attention to the fact that there is an
infinite difference between the interest of faith and the interest of
knowledge: His Rev. understands it as if one were speaking
to him of ``an indifferent and disinterested knowledge.''

One draws Dr. M.'s attention to how the
speculation ``he has acknowledged in his Dogmatics''
% --- p. 64 ---
\setcounter{footnote}{0}%
\opage{64}dissolves through its own inner contradiction: His Rev. understands
it as if one took no account at all of ``the speculation''
he has acknowledged in his Dogmatics.

One draws Dr. M.'s attention to the fact that a systematic
head (be it philosopher or theologian) who wants to assert
``coherent cognition'' and chastise those who must
make do with ``stray thoughts, aphorisms, fancies and flashes''
must above all apply himself to having a principle and respecting
its consequences: His Rev. understands it as if
one thereby committed a --- ``\textit{metabasis in aliud genus}!''

One draws Dr. M.'s attention to the fact that he cannot possibly
bring speculation along if he really grants
that the thinking of faith is ``an existential thinking.'' His Rev.
understands it as if one had forgotten that he too in his
Dogmatics grants that the thinking of faith is existential.

One draws Dr. M.'s attention to the fact that faith certainly
has its object, its definite object, but that the cognition of truth
in faith relates itself to this object precisely
in such a way that subjective certainty, subjective appropriation,
here becomes the one decisive thing: His Rev. understands it as
if one were speaking to him of a cognition of truth that
had nothing objective, and instructs us (Oplysn. p. 16) that
``even the poorest rationalist dogmatics'' must after all recognize
``that the truth has a certain objectivity.''

One draws Dr. M.'s attention to the fact that, since the object of
faith must necessarily correspond to faith and the cognition of faith,
the object of faith for precisely this reason
cannot be an object of knowledge, since faith as faith is precisely
the opposite of objective knowledge: His Rev. understands it
% --- p. 65 ---
\setcounter{footnote}{0}%
\opage{65}(cf. Oplysn. pp. 19—22) as if one were speaking to him of
a faith that was supposed to be in conflict with its object.

Finally, one draws Dr. M.'s attention to the fact that the objective
incomprehensibility of the objects of faith precisely becomes, and must
become, evident to him who knows what it means to comprehend
something, and who besides thinks earnestly and persistently about
his faith. Dr. M. understands it as if one had now grown
tired of thinking about one's faith, and exclaims with Anselm:
``\textit{negligentiæ mihi videtur},\dots{} Yes: \textit{negligentiæ mihi videtur}!''
---

Such, then, in the main is Dr. M.'s understanding of the thoughts
of his ``opponents.'' Indeed, what is one to say? If
such an understanding is now also to be a sign of the ``redemption
of thought,'' then it looks --- at least to my eyes ---
somewhat motley with the ``logic of rebirth.''

About the attempts which the author of ``dogm. Opl.'' has here and there
made at explaining and supporting his so ``misinterpreted''
concept of science, I shall now be as brief as possible.
a) ``Dogmatic speculation'' cannot be defended
by an appeal to the ``consciousness of revelation''; for
the facts of revelation are --- as we have seen
--- only for the existential appropriation of faith: a speculative
consciousness of revelation is a contradiction.

b) ``Dogmatic speculation'' cannot be defended
by reference to the ``contemplation'' found in the Psalms
of David and in the Epistles of Paul; to speculate is not the
same as to sing psalms; to convert heathens is
not the same as to speculate on the Trinity: a
scientific psalmist, a speculating apostle, is a
contradiction.

% --- p. 66 ---
\setcounter{footnote}{0}%
\opage{66}c) ``Dogmatic speculation'' cannot be defended
by reference to the old Fathers, Tertullian, Augustine,
Anselm etc.: the dear Fathers could well be good
Christians and yet find themselves in some unclarity about the relation
between ``\textit{certitudo experientiæ}'' and ``\textit{certitudo}
\textit{speculativa}'', just as the dear Dr. Martensen too,
who still, despite everything that the scientific thinking
of the centuries has elucidated, seems to want to petrify himself in the same
unclarity, can nonetheless for his own person well be a
good Christian.

d) Finally, ``dogmatic speculation'' cannot be defended
either by reference to the mutual coherence of the doctrines
of faith. Only speculative propositions are in
speculative coherence. Were the coherence between
the dogmas to be speculative, then each single dogma
(thus also the dogma of the omnipresence of the transfigured Christ)
would, after all, have to be speculative; but with this the case
is now, as we have seen, otherwise: a speculative confession of faith,
a speculative catechism, a speculative dogmatics are
contradictions which not even the ``logic of rebirth'' shall be able
to ``mediate.''

Still, as I said, I am being brief, partly because I now
assume that at least the point of contention must be visible to
those who can and will see, partly also because I cannot
conclude this discussion without once more directing
the reader's attention to Dr. M.'s and my position in relation to
``the pseudonym'' and the Kierkegaardian literature.

Instead of keeping to my review of ``the Christian
Dogmatics,'' the author of dogm. Opl. has preferred
to keep to my excerpt, i.e. to my review of
% --- p. 67 ---
\setcounter{footnote}{0}%
\opage{67}``Johannes Climacus.'' --- By thus striking together two magnitudes
as heterogeneous as ``the pseudonym'' and my humble self
into one, His Rev. has no doubt intended to strike,
as the saying goes, ``two flies with one swat,'' since he
surely thought he could strike me in ``the pseudonym'' and
then at the same stroke strike ``the pseudonym'' in me.

Has Dr. M. now hit me in ``the pseudonym''? --- I
cannot perceive it. For as regards my having
spoken in my ``master's words,'' I could, after all, not well do
otherwise when I wanted to give a reliable ``excerpt'' of his
``doctrines.'' As regards my not even having
tried to explain to His Rev. what I mean by ``\emph{the paradox}'':
this is only to be understood as meaning that His Rev.
has not even tried to think through what stands expressly
in my writings. As regards, finally, my supposedly
having founded a standpoint out of ``frivolity'' and ``dilettantish inconstancy,''
merely in order to belong among the ``exceedingly
interesting,'' it would surely first have to be demonstrated that my
chosen ``standpoint'' is untenable. The intended ``swat''
therefore cannot hit me before it has hit ``the pseudonym.''


Has Dr. M., then, really hit ``the pseudonym'' in me?
--- I will not insult Mr. Kierkegaard and make myself
ridiculous in the eyes of the ``pseudonyms'' by defending ``Johannes
Climacus'' against Dr. Martensen; I will only with a couple of words
illuminate Dr. M.'s opinion of what he calls the pseudonyms'
``prolix literature'' by referring to the remarks tossed out
in his ``Dogmatic Elucidations.''

For a number of years now the pseudonymous thinkers,
one after the other, have fixed attention on
% --- p. 68 ---
\setcounter{footnote}{0}%
\opage{68}the tasks of ``inwardness'' and summoned all their dialectic to
make the Christian problem recognizable; for a number
of years Dr. M., who as dogmatician occupies himself with
the same problem, has been aware that these pseudonyms,
although far too inward-turned to express themselves with any
finite polemic of personalities, nevertheless constantly kept
a watchful eye on ``speculation'' and ``the system''; for
a number of years a quietly growing attention has followed
this so wonderfully repellent and yet so infinitely attractive
``literature,'' so that many must surely many times
have asked themselves: what is to be done about this? --- Under
these circumstances we must, then, be doubly grateful
to Dr. M. that he finally steps forward, and not only tells us
what is to be done about this, but also tells it to us in so
``direct a communication.'' From this ``direct communication''
everyone must surely be able to see that Dr. M. has not read, and
surely does not intend to read either, ``the writings of the pseudonyms,'' that
His Rev. cannot understand these writings\footnote[1]{Dr. Martensen's definite declaration that he cannot understand ``the pseudonyms'' might, however, here possibly be taken in a more literal sense than His Rev. really wishes. What he thus, e.g. (Opl. p. 39), surely not so much after Anselm's as after Frederikke Bremer's ``prototype,'' explains to us about the pillar-hero, ``the single individual,'' item the ingenious reflections he makes here and there on the ``passion'' of faith and the ``risk'' of faith, seem really to indicate that his ``dogmatic genius'' is in this respect legitimately excused. What here in fairness speaks for Dr. M.'s excuse is, then, also this, that the pseudonyms offer him absolutely no fixed and reliable ποῦ στῶ.\par The dialectical pseudonyms have, as it seems, all arranged themselves according to the Lichtenbergian motto of the ``Stages'': ``Solche Werke sind Spiegel; wenn ein Affe hinein guckt, kann kein Apostel heraus sehen.'' Such a ``standpoint'' is undeniably only a little ``dogmatic''; but --- one must take it as it is, since no other ποῦ στῶ is offered.}, ``that in accordance with the course of his
% --- p. 69 ---
\setcounter{footnote}{0}%
\opage{69}studies'' and his ``individual intellectual orientation'' he does not exactly
wish to understand these writings either, and that therefore the
whole cannot mean much either.

It is a hard blow, indeed, it is a deadly blow,
the blow with which Dr. M.'s dogmatic superiority here
has struck the only slightly dogmatic ``Johannes Climacus'';
but --- this I must confess --- although this blow, as regards
reality, is so hard a blow, Dr. M. has, as regards
form, nevertheless known how to ``observe the limit'' and
deliver the blow with all the composure and apparent forbearance
that always becomes the superior one.

Dr. M. declares, to be sure, that he has no use whatever
for the ``doctrines'' which I have wished to ``impart'' to him
from ``the pseudonym,'' since he must apply to these the old
saying ``that what is true in them is not new, and what is new not true'';
but he nevertheless also grants that ``the pseudonym'' may to
a certain degree be excused by the fact that, as a ``gymnast''
and ``humorist,'' he merely intends ``in Socratic
fashion to set a deeper skepticism in motion (Oplysn.
p. 12).'' Most of the blame thus falls on me, who have
been so ``frivolous'' as to take the humorist's ``propositions'' for
sheer earnest, and to think that these propositions could (Oplysn.
p. 13) be used for ``a new doctrine of cognition and a treatment
of theology grounded on it.'' --- But, as I said,
this forbearance toward the ``gymnast'' lies solely in
the form. The dogmatic verdict runs thus:

If what ``Johannes Climacus'' says about the Christian
problem is to be understood as a gymnastic-humorous half-jest,
there may perhaps well be something in it; but if his
main propositions are to be taken in earnest and to count as
% --- p. 70 ---
\setcounter{footnote}{0}%
\opage{70}truth, then Dr. M. must most definitely declare that ``what is
true in them is not new, and what is new not true''; and so in reality
it is, after all, the pseudonym who pays for it nonetheless.


As for my opinion, I say to all this
quite directly: read the Kierkegaardian writings! Reader, you
who want to have some judgment in this matter, whoever you are,
learned or unlearned, just read the Kierkegaardian writings in all
quietness, without being disturbed by the buzzing of dogma and the strife
of parties. Read these writings with intellectual interest!
test what you read by yourself alone, and you shall see
how lightly and playfully, and yet --- how earnestly, the ``humorous''
authors unite in leading you into ``the problem,'' into
the true problem of inwardness, of personality, of eternity,
into the problem for whose sake alone it is, after all, worth
living. Or: read these writings with religious-Christian
interest! free your mind from all the outward-turned,
worldly-ecclesiastical passions of finitude, apply what you have read to
yourself, and you shall see with what instinctive and yet conscious
certainty these otherwise so different authors know how to show
you the Christian in Christianity, and what mildness, reconciliation
and peace nevertheless here hover over the many otherwise so
disquieting and deeply terrifying thoughts. Or, if you
would rather have it so: read these writings, then, in decidedly polemical
interest! Assume that what I have said here is merely a
foolish exaggeration, sprung either from blind enthusiasm
or from crafty flattery of another human being whose recognition
is perhaps expected in return, if only in lesser
measure; assume that the whole is nothing but an arbitrarily
defiant ``one-sidedness'' which must most emphatically be
% --- p. 71 ---
\setcounter{footnote}{0}%
\opage{71}refuted; resolve earnestly that you will be precisely the one who
comes forward with the refutation, choose your definite points, make
a draft, but --- meanwhile read all the writings concerned,
``the pseudonyms'' as well as the ``edifying discourses,'' read
them thoroughly (for if you do not read thoroughly, you surely cannot
refute thoroughly either): you shall see that it does not, after all,
go quite so easily with the refutation; it might
easily happen to you that, by entering into the actual
context, you would here find a self-criticism that made your
objections superfluous, and in the end impounded your whole intended
critique.

That I, for the rest, am not in possession of any direct
or authentic explanation of these writings is self-evident,
since every direct understanding is here made dialectically impossible.
That the impossibility of such a direct understanding
should serve as proof that ``indirect communication''
is imperfect and of a lower kind than direct is, however,
an erroneous opinion that is easy to correct.

It has, to be sure, been thought that ``the pseudonyms,'' with
their humorously playful form and their doubly reflected
categories, must be less earnest and therefore less
reliable than the ``edifying discourses'' (e.g. ``Works
of Love''). The opposite is precisely the case. Precisely by % print: Gjerninger„ (misprint)
their indirect form of communication these ``pseudonyms'' stand,
in earnest, on an equal footing with the edifying writings\footnote[1]{Whose communication, moreover, although in an entirely different sense, is and must be indirect.},
(indeed, are, with all their wit and humor, at bottom far stricter);
by their ``indirect form of communication'' the
% --- p. 72 ---
\setcounter{footnote}{0}%
\opage{72}pseudonyms achieve precisely what is absolutely necessary when one
is to speak absolutely about the absolute: they achieve, namely, by
their transformations in the comic, the cutting off of all relative
reasoning and the halting of the otherwise unhaltable discussion.


Had, e.g., ``Johannes Climacus'' delivered his ``doctrines''
reviewed by me in a pure and direct scientific
form, it would have become a serious matter for the theologians,
and they would then surely also have taken it very
seriously. Now, on the other hand, since he has communicated these doctrines
in the indirect form of wit and humor, now one does not regard
the matter as so serious. Herein is seen a confusion
of the higher and the lower.

With a merely scientific form ``the author''
would evidently have moved in a circle, insofar as he,
although polemicizing against science, nevertheless remained
within science, and thus attempted, in that contradiction,
to make the unscientific scientific. With
a merely scientific form ``the author'' would have found
himself in the contradiction of speaking out of the ground of subjectivity and
inwardness without, however, speaking the language of subjectivity and
inwardness, of saying ``directly'' what no
human being can, after all, see or understand ``directly.'' When, then, an author
who wants to carry a principle through and make a problem
recognizable uses precisely the form of communication which, by
the nature of the matter, is the only one adequate to his problem,
shall we not then agree to admit that the choice
of such a form indicates precisely that this author has
--- earnestness; and when it is finally seen that the problem of inwardness
can only be communicated in the form of inwardness, and
% --- p. 73 ---
\setcounter{footnote}{0}%
\opage{73}that this form, where a position over against
``science'' is concerned, must necessarily become a form which, as
it unfolds the content of subjectivity, at the same time secures this
content by every utterance humorously canceling its
immediate objective validity, in order to resound all the more strongly % print: før (misprint)
in inwardness itself and rest in subjectivity: shall we
not then agree to admit that far
richer gifts of spirit and a far higher productivity are required to % print: før (misprint)
communicate the problem of inwardness in the language of inwardness
than are needed to set forth one's thoughts in a direct
scientific form.

It will easily be seen from what has been adduced here that I regard
the Kierkegaardian literature as something in itself so
peculiar, self-contained and unique that there can be no thought
at all here of getting these writings introduced by any direct
way and incorporated into the rest of the literature. But
from this it by no means follows that one should be permitted in ``literature''
to ignore these writings or act as if one
had not read them. On the contrary! Anyone who nowadays
wants to join in talking about the relation between faith and cognition
(be he theologian or philosopher, clergyman or layman) will
surely not clear up much for us if he himself is
in unclarity about the ``problem'' of these writings.

Insofar as the Kierkegaardian writings belong to the public,
they naturally also belong to literature. Every
author whose task touches theirs can and ought to place himself
in an entirely definite and unambiguous relation to them.

Whether this relation is to be polemically dismissive or benevolently
appreciative, about that let each ask himself, without
taking Mr. Kierkegaard into counsel. Should Mr.
% --- p. 74 ---
\setcounter{footnote}{0}%
\opage{74}Kierkegaard then have something to object, he can, after all, speak up, and himself
join in the discussion.

That I for my part have not only personal goodwill toward
Mr. Kierkegaard but also great esteem for the man
from whose writings I can see clearly enough (if not ``directly'')
that his inner life must be far more strenuous
than mine, I surely need not assure further. But
it would grieve me deeply if Mr. Kierkegaard
(which I must surely consider unthinkable) should misunderstand
these remarks of mine, as if I thereby wished to place
myself in any relation of dependence whatsoever on his personal individuality. My relation to this authorial individuality
is, and must essentially be, the same as it
would be if he, instead of living here in Copenhagen,
lived e.g. in Rome or in Philadelphia, or as if he,
instead of being our ``highly esteemed contemporary,'' were assumed to have
lived either with Tertullian or with Augustine or with
--- Anselm.

To come into an untrue relation to the Kierkegaardian
writings I consider extremely dangerous, insofar as I always
consider it extremely dangerous to come into an untrue relation
to a truth; to come into a literary polemic with Mr.
Kierkegaard I do not, on the other hand, in this sense consider
so dangerous. Let him have ever so much shrewdness and ever so
brilliant weapons; good Lord, we cannot, after all, all be geniuses;
when one fights for a cause, one does not, after all, fight in order
to show one's talent: Mr. Kierkegaard is, moreover, as far as I
know, precisely the only author who has summoned, indeed I could
say: sacrificed and squandered, all his talent on proving that for
the cognition of the highest truth talent is a matter of
% --- p. 75 ---
\setcounter{footnote}{0}%
\opage{75}indifference. Let Mr. Kierkegaard show that there is either
ambiguity, unclarity or other impropriety in the use
I have made of his writings: if he can convince me,
well then, I confess my blunder and thank him for his
guidance; if he cannot convince me, well, then I stick
to my own, and should think that I have in this respect the same
right as anyone else.

Thus: my cause is not Mr. Kierkegaard's cause.
Mr. Kierkegaard's cause is outside all discussion, ideally
fixed in that center of inwardness about which the religious
moves. My cause is in the discussion, turned outward
toward the boundary line where theology and philosophy
contend about knowledge and non-knowledge.

When I published my first writing, ``The Faith of the Gospel and
the Modern Consciousness,'' I thought I could conduct my
cause against theology without touching any of our theologians
personally. Dr. Martensen's attitude in the preface to
his Dogmatics prompted me to review this book
in such a tone that his genteel ambiguity would necessarily
have to burst. Meanwhile one theologian after another has felt
offended; but although in my
second writing, ``The Faith of the Gospel and Theology,'' I have given a detailed
account of my conception of the relation of Christianity
to science, no theologian has yet attempted either
to overturn my ``propositions'' or at all to acquaint himself
``scientifically'' with the matter.

\end{document}
